\documentclass[11pt]{article}

\usepackage[margin=1in]{geometry}
\usepackage{amsmath,amssymb,mathtools}
\usepackage{booktabs,longtable,array}
\usepackage[T1]{fontenc}
\usepackage{lmodern}
\usepackage{textcomp}
\usepackage[hidelinks]{hyperref}
\numberwithin{equation}{section}
\allowdisplaybreaks

\newcommand{\EM}{\textsc{Emmanuel}}
\newcommand{\NS}{\textsc{Newspeak}}
\newcommand{\kw}[1]{\ensuremath{\mathsf{#1}}}
\newcommand{\str}[1]{\ensuremath{\text{\texttt{#1}}}}
\newcommand{\bs}{\textbackslash}
\newcommand{\Open}{\str{[\bs}}
\newcommand{\Close}{\str{\bs]}}
\newcommand{\TM}{\kw{T}}
\newcommand{\CM}{\kw{C}}
\newcommand{\lexto}{\longrightarrow_{\mathrm{lex}}}
\newcommand{\rulelabel}[1]{\textsc{#1}}
\newcommand{\vs}{\text{\textvisiblespace}}
\newcommand{\crt}{\mathord{\curlywedge}}

\title{A Formal Semantics for \EM\\\large Design draft 0.1: the parametric core}
\author{}
\date{\today}

\begin{document}
\maketitle

\begin{abstract}
\EM{} is the mark-edit language of Ampleforth, Newspeak's live document
system.  A mark-edit language is a markup language designed so that a
document can be edited both as source and through its live rendering, with
every edit made in the rendering carried back to source.  This draft gives
the parametric core of \EM: its lexical structure, its grammar, the
interpretation of markup into a tree of document blocks, and printing.  The
core is parametric in three things: the host that parses and runs the
\NS{} code occurring in markup, the block signature that says what
each message does to each kind of block, and the field interface used by
a live view.  On the core it builds the structural core of the standard
library and the edit semantics: how each edit made in the live view is
carried back to source as a rewrite of it.
\end{abstract}

\clearpage
\begingroup
\footnotesize
\tableofcontents
\endgroup
\clearpage

\section{Authority, scope, and method}
\label{sec:method}

The normative source is the Ampleforth document
\emph{Newsepak\_Markup\_Design\_Notes}, in the state of its decisions dated
November 7, 2026 (henceforth \emph{the design notes}).  Where this document
and the design notes disagree, the disagreement is a defect to be resolved
with the author; neither silently governs.

The specification is layered:
\begin{enumerate}
\item lexical structure (Section~\ref{sec:lexical});
\item grammar and abstract syntax;
\item interpretation, a small-step transition system over a block store,
      a stack of open blocks, and the sequence of parsed document blocks;
\item printing, with the round-trip property stated as an obligation;
\item the standard library, the block signature of ordinary documents
      (Section~\ref{sec:library});
\item edit semantics: provenance, and live-view edits as
      source rewrites.
\end{enumerate}

\paragraph{Parameters.}
The core does not embed \NS.  It is parametric in
\begin{itemize}
\item a \emph{host}, with three components: a \emph{host parser}, which
      parses the code of a bracket into statements
      (Section~\ref{sec:host-parser}), a \emph{host evaluator}, which
      evaluates a \NS{} expression in the scope of the document and may
      diverge, and a \emph{host printer}, which prints statements back
      (Section~\ref{sec:printing});
\item a \emph{block signature}, which for each kind of block and message
      gives the message's effect: pop, add a leaf, add and push, mutate,
      refuse, pop and reject, or not understood; and
\item a \emph{field interface}, used only by the edit layer
      (Section~\ref{sec:field}).
\end{itemize}
The standard block library (sections, paragraphs, styles, tables, citations,
and so on) is one instance of the block signature; its structural core is
given in Section~\ref{sec:library}.  Ordinary \NS{} supplies the host through the
\NS{} semantics.

\paragraph{Conventions.}
As in the \NS{} semantics, disjunctive conditions are written as separately
named rules with no first-match priority, chains and closures are inductive
relations without fuel, and absent results have explicit constructors.

%----------------------------------------------------------------------
\section{Lexical structure}
\label{sec:lexical}

\subsection{Characters, positions, and spans}

Let $\Sigma$ be the set of Unicode scalar values and let $w \in \Sigma^*$
range over character sequences; $\epsilon$ is the empty sequence and
juxtaposition is concatenation.  A markup source is a sequence
$W = c_0 c_1 \cdots c_{n-1}$.  A \emph{position} is an index
$0 \le i \le n$, and a \emph{span} $[i,j)$ with $i \le j$ denotes
$c_i \cdots c_{j-1}$.  Every token and every node of abstract syntax carries
the span of the source it was read from; the design notes require this so
that a block can recover the exact source of every message it received.

Four two-character sequences are distinguished:
\begin{equation}\label{eq:lexical-sequences}
\Open = \str{[}\,\str{\bs} \qquad
\Close = \str{\bs}\,\str{]} \qquad
\str{\bs\bs} \qquad
\str{\bs[}
\end{equation}
The first two are \emph{mode switches}; the last two are the only
\emph{escapes}, and occur only in text.

\subsection{Modes}

The lexer has two modes, text $\TM$ and code $\CM$, and keeps a stack of
them.  Let $\mu$ range over mode stacks, written with the top at the right:
\begin{equation}\label{eq:mode-stacks}
\mu ::= \TM \mid \mu\,\TM \mid \mu\,\CM
\end{equation}
The initial stack is $\TM$.  The bottom of every stack is $\TM$, and $\TM$
and $\CM$ alternate: $\Open$ always pushes the mode other than the current
one, and $\Close$ pops.  Hence a stack is determined by its depth, and
\begin{equation}\label{eq:mode-alternation}
\mu = \TM\,\CM\,\TM\,\CM \cdots
\end{equation}
A $\CM$ pushed over $\TM$ opens a bracket of messages; a $\TM$ pushed over
$\CM$ opens a markup expression nested in code.

\subsection{Lexemes}

The lexer produces a sequence of \emph{lexemes}, each with its span:
\begin{equation}\label{eq:lexemes}
\lambda ::= \kw{char}(c) \mid \kw{openMessages} \mid \kw{closeMessages}
  \mid \kw{openMarkup} \mid \kw{closeMarkup} \mid \kw{code}(w)
\end{equation}
Here $\kw{char}(c)$ is one character of text after escapes are resolved,
and $\kw{code}(w)$ is one piece of code source: a single character, a whole
string literal, or a whole comment.

\subsection{The lexing relation}

A lexer configuration is $\langle \mu, w \rangle$, a mode stack and the
input not yet read.  A step reads a prefix of $w$ and emits one lexeme:
$\langle \mu, w \rangle \lexto^{\lambda} \langle \mu', w' \rangle$.
Spans are omitted from the rules; the span of each lexeme is the prefix it
consumes.  In the rules, $\mu$ ranges over possibly empty prefixes of a mode
stack, so that $\mu\TM$ matches every stack whose top is $\TM$, including the
initial stack.

\paragraph{Text mode.}
\begin{gather}
\label{eq:lex-open-messages}
\rulelabel{Lex-Open-Messages}\quad
\langle \mu\TM, \Open\,w \rangle \lexto^{\kw{openMessages}}
  \langle \mu\TM\CM, w \rangle
\\
\label{eq:lex-close-markup}
\rulelabel{Lex-Close-Markup}\quad
\langle \mu\CM\TM, \Close\,w \rangle \lexto^{\kw{closeMarkup}}
  \langle \mu\CM, w \rangle
\\
\label{eq:lex-escape-backslash}
\rulelabel{Lex-Escape-Backslash}\quad
\langle \mu\TM, \str{\bs\bs}\,w \rangle \lexto^{\kw{char}(\str{\bs})}
  \langle \mu\TM, w \rangle
\\
\label{eq:lex-escape-bracket}
\rulelabel{Lex-Escape-Bracket}\quad
\langle \mu\TM, \str{\bs[}\,w \rangle \lexto^{\kw{char}(\str{[})}
  \langle \mu\TM, w \rangle
\\
\label{eq:lex-text-char}
\rulelabel{Lex-Text-Char}\quad
\langle \mu\TM, c\,w \rangle \lexto^{\kw{char}(c)}
  \langle \mu\TM, w \rangle
  \quad\text{if } c\,w \text{ begins with none of }
  \Open, \Close, \str{\bs\bs}, \str{\bs[}
\end{gather}
$\Close$ in text mode with stack $\TM$, that is, outside any markup
expression, has no rule; it is the lexical error
$\kw{strayClose}$ (Section~\ref{sec:lexical-errors}).

\paragraph{Code mode.}
\begin{gather}
\label{eq:lex-open-markup}
\rulelabel{Lex-Open-Markup}\quad
\langle \mu\CM, \Open\,w \rangle \lexto^{\kw{openMarkup}}
  \langle \mu\CM\TM, w \rangle
\\
\label{eq:lex-close-messages}
\rulelabel{Lex-Close-Messages}\quad
\langle \mu\TM\CM, \Close\,w \rangle \lexto^{\kw{closeMessages}}
  \langle \mu\TM, w \rangle
\\
\label{eq:lex-code-string}
\rulelabel{Lex-Code-String}\quad
\langle \mu\CM, \str{'}\,w \rangle \lexto^{\kw{code}(\str{'}\,u)}
  \langle \mu\CM, w' \rangle
  \quad\text{if } \mathit{lit}(w) = \langle u, w' \rangle
\\
\label{eq:lex-code-comment}
\rulelabel{Lex-Code-Comment}\quad
\langle \mu\CM, \str{(*}\,w \rangle \lexto^{\kw{code}(\str{(*}\,u)}
  \langle \mu\CM, w' \rangle
  \quad\text{if } \mathit{com}(1, w) = \langle u, w' \rangle
\\
\label{eq:lex-code-char}
\rulelabel{Lex-Code-Char}\quad
\langle \mu\CM, c\,w \rangle \lexto^{\kw{code}(c)}
  \langle \mu\CM, w \rangle
\notag\\
  \text{if } c \ne \str{'} \text{ and } c\,w \text{ begins with none of }
  \Open, \Close, \str{(*}
\end{gather}
A string literal and a comment are each read whole by a \emph{reader}, a
function from the input after the opening delimiter to either a pair
$\langle u, w' \rangle$, the rest $u$ of the literal (closing delimiter
included) and the input $w'$ after it, or $\bot$ if the input ends first.
Write $v \cdot \langle u, w' \rangle = \langle v\,u, w' \rangle$ and
$v \cdot \bot = \bot$.  The string reader follows the \NS{} string literal:
a quote is doubled inside the literal, and the literal ends at the first
quote that is not doubled.
\begin{equation}\label{eq:string-reader}
\begin{aligned}
\mathit{lit}(\str{'}\,\str{'}\,w) &= \str{'}\,\str{'} \cdot \mathit{lit}(w) \\
\mathit{lit}(\str{'}\,w) &= \langle \str{'}, w \rangle
  && \text{if } w \text{ does not begin with } \str{'} \\
\mathit{lit}(c\,w) &= c \cdot \mathit{lit}(w) && \text{if } c \ne \str{'} \\
\mathit{lit}(\epsilon) &= \bot
\end{aligned}
\end{equation}
The comment reader follows the \NS{} comment, which nests; its first
argument $d \ge 1$ is the number of comments open.
\begin{equation}\label{eq:comment-reader}
\begin{aligned}
\mathit{com}(1, \str{*)}\,w) &= \langle \str{*)}, w \rangle \\
\mathit{com}(d, \str{*)}\,w) &= \str{*)} \cdot \mathit{com}(d-1, w)
  && \text{if } d > 1 \\
\mathit{com}(d, \str{(*}\,w) &= \str{(*} \cdot \mathit{com}(d+1, w) \\
\mathit{com}(d, c\,w) &= c \cdot \mathit{com}(d, w)
  && \text{if } c\,w \text{ begins with neither } \str{(*} \text{ nor } \str{*)} \\
\mathit{com}(d, \epsilon) &= \bot
\end{aligned}
\end{equation}
The equations of each reader have disjoint conditions, so each reader is a
function.  Because a string literal or comment is read whole, a mode switch
inside one has no effect; this is how the design notes' lexical rule~4 is
realized.  A string literal or comment whose reader yields $\bot$ is the
lexical error $\kw{unclosedLiteral}$.

Note that in code mode $\str{\bs\bs}$ and $\str{\bs[}$ are not escapes:
$\str{\bs\bs}$ is the \NS{} modulo operator, read as two
$\kw{code}$ lexemes.

\subsection{Determinism}

\begin{equation}\label{eq:lex-determinism}
\langle \mu, w \rangle \lexto^{\lambda_1} \langle \mu_1, w_1 \rangle
\ \wedge\
\langle \mu, w \rangle \lexto^{\lambda_2} \langle \mu_2, w_2 \rangle
\implies \lambda_1 = \lambda_2 \wedge \mu_1 = \mu_2 \wedge w_1 = w_2
\end{equation}
In text mode the four distinguished sequences differ in their first two
characters, so at most one of the first four rules applies, and
\rulelabel{Lex-Text-Char} applies exactly when none does.  In code mode,
$\Open$, $\Close$, a quote, and $\str{(*}$ are mutually exclusive
beginnings, and the readers are functions.  Lexing is therefore a
deterministic left-to-right transduction.  For example, the text
$\str{\bs\bs]}$ reads as $\kw{char}(\str{\bs})\,\kw{char}(\str{]})$, not as
$\str{\bs}$ followed by $\Close$, because the escape begins first.  This is
why the design notes write the literal text $\Close$ as $\str{\bs\bs]}$.

\subsection{Lexing a source}
\label{sec:lexical-errors}

A \emph{located lexeme} $\ell = (\lambda, [i,j))$ pairs a lexeme with its
span.  Running the lexer from position $i$ yields a result $r$:
\begin{equation}\label{eq:lex-results}
r ::= \kw{lexed}(\bar\ell) \mid \kw{error}(e, i)
\qquad
e ::= \kw{strayClose} \mid \kw{unclosedLiteral} \mid \kw{unclosedBracket}
\end{equation}
The run judgment $\langle \mu, i, w \rangle \Downarrow r$ says that lexing
$w$, which begins at position $i$, from mode stack $\mu$ yields $r$.
\begin{gather}
\label{eq:run-end}
\rulelabel{Run-End}\quad
\langle \TM, i, \epsilon \rangle \Downarrow \kw{lexed}(\epsilon)
\\
\label{eq:run-step}
\rulelabel{Run-Step}\quad
\frac{\langle \mu, w \rangle \lexto^{\lambda} \langle \mu', w' \rangle
      \qquad j = i + |w| - |w'|
      \qquad \langle \mu', j, w' \rangle \Downarrow \kw{lexed}(\bar\ell)}
     {\langle \mu, i, w \rangle \Downarrow
      \kw{lexed}((\lambda, [i,j))\,\bar\ell)}
\\
\label{eq:run-step-error}
\rulelabel{Run-Step-Error}\quad
\frac{\langle \mu, w \rangle \lexto^{\lambda} \langle \mu', w' \rangle
      \qquad j = i + |w| - |w'|
      \qquad \langle \mu', j, w' \rangle \Downarrow \kw{error}(e, i')}
     {\langle \mu, i, w \rangle \Downarrow \kw{error}(e, i')}
\\
\label{eq:run-stray-close}
\rulelabel{Run-Stray-Close}\quad
\langle \TM, i, \Close\,w \rangle \Downarrow \kw{error}(\kw{strayClose}, i)
\\
\label{eq:run-unclosed-string}
\rulelabel{Run-Unclosed-String}\quad
\langle \mu\CM, i, \str{'}\,w \rangle \Downarrow
  \kw{error}(\kw{unclosedLiteral}, i)
  \quad\text{if } \mathit{lit}(w) = \bot
\\
\label{eq:run-unclosed-comment}
\rulelabel{Run-Unclosed-Comment}\quad
\langle \mu\CM, i, \str{(*}\,w \rangle \Downarrow
  \kw{error}(\kw{unclosedLiteral}, i)
  \quad\text{if } \mathit{com}(1, w) = \bot
\\
\label{eq:run-unclosed-bracket}
\rulelabel{Run-Unclosed-Bracket}\quad
\langle \mu, i, \epsilon \rangle \Downarrow \kw{error}(\kw{unclosedBracket}, i)
  \quad\text{if } \mu \ne \TM
\end{gather}
The span of each lexeme is the prefix of the input its step consumes.  An
error records the position where it is detected: the stray $\Close$, the
start of the unclosed literal, or the end of the input.  A
source $W$ \emph{lexes} to $\bar\ell$, or \emph{fails} with $e$ at $i$, according
to the run from the initial configuration:
\begin{equation}\label{eq:lex-source}
W \text{ lexes to } \bar\ell \iff
  \langle \TM, 0, W \rangle \Downarrow \kw{lexed}(\bar\ell)
\qquad
W \text{ fails with } e \text{ at } i \iff
  \langle \TM, 0, W \rangle \Downarrow \kw{error}(e, i)
\end{equation}
On the stacks reachable from $\TM$, which alternate
by~\eqref{eq:mode-alternation}, every configuration either takes a lexing
step or matches exactly one of the error rules, and none does both: a
$\Close$ in text mode can step only when the stack has a $\CM$ under its
top, and a quote or $\str{(*}$ in code mode steps exactly when its reader
does not yield $\bot$.  With determinism~\eqref{eq:lex-determinism}, every
source therefore has exactly one result.

\paragraph{Obligation (no collision with \NS).}
The design notes' lexical rule~4 asserts that neither $\Open$ nor $\Close$
can occur in valid \NS{} code outside string literals and comments.  This is
a property of the \NS{} grammar, recorded here as an obligation to be
discharged against the \NS{} front-end semantics.  It guarantees that the
lexer never misreads valid code as a mode switch.

\subsection{Tokens}

The parser consumes \emph{tokens}, obtained from located lexemes by
coalescing.
\begin{equation}\label{eq:tokens}
\tau ::= \kw{text}(u) \mid \kw{code}(w) \mid \kw{openMessages}
  \mid \kw{closeMessages} \mid \kw{openMarkup} \mid \kw{closeMarkup}
\end{equation}
A located token $t = (\tau, [i,j))$ pairs a token with its span.  Each
lexeme first becomes a token, $\kw{char}(c)$ becoming $\kw{text}(c)$ and
every other lexeme the token of the same name; then every maximal run of
$\kw{text}$ tokens is merged into one, and likewise every maximal run of
$\kw{code}$ tokens.
\begin{equation}\label{eq:coalesce}
\begin{aligned}
\mathit{tokens}(\bar\ell) &= \mathit{merge}(\mathit{tok}(\ell_1) \cdots
  \mathit{tok}(\ell_k)) \qquad \text{where } \bar\ell = \ell_1 \cdots \ell_k \\
\mathit{merge}(\epsilon) &= \epsilon \\
\mathit{merge}(t) &= t \\
\mathit{merge}((\kw{text}(u), [i,j))\,(\kw{text}(u'), [j,k))\,\bar t)
  &= \mathit{merge}((\kw{text}(u\,u'), [i,k))\,\bar t) \\
\mathit{merge}((\kw{code}(w), [i,j))\,(\kw{code}(w'), [j,k))\,\bar t)
  &= \mathit{merge}((\kw{code}(w\,w'), [i,k))\,\bar t) \\
\mathit{merge}(t\,t'\,\bar t) &= t\,\mathit{merge}(t'\,\bar t)
\\
  &\phantom{{}={}} \text{if } t, t' \text{ are neither both } \kw{text}
     \text{ nor both } \kw{code}
\end{aligned}
\end{equation}
Consecutive lexemes of a run have adjacent spans, so the merging equations
always apply to them, and the span of a merged token is the union of the
spans merged.  Coalescing realizes the maximal-munch reading of text on
which the design notes' argument for the trailing $\kw{Text}$ depends: two
$\kw{text}$ tokens are never adjacent.

%----------------------------------------------------------------------
\section{Grammar and abstract syntax}
\label{sec:grammar}

The design notes give the grammar
\begin{equation}\label{eq:concrete-grammar}
\begin{aligned}
\kw{DocumentMarkup} &\to \kw{DocumentBlock}^{*}\ \kw{Text} \\
\kw{DocumentBlock} &\to \kw{Text}\ \Open\ \kw{Messages}\ \Close \\
\kw{MarkupExpression} &\to \Open\ \kw{DocumentMarkup}\ \Close
\end{aligned}
\end{equation}
where $\kw{Text}$ may be empty, $\kw{Messages}$ is a sequence of \NS{}
statements separated by dots, each a single message without an explicit
receiver and without cascades, and a $\kw{MarkupExpression}$ may occur in
$\kw{Messages}$ wherever \NS{} allows an expression.

Parsing is split in two.  The \emph{structural parse} reads the token
sequence of Section~\ref{sec:lexical} into raw nodes, in which the code of a
bracket is kept as a sequence of pieces: code, and the markup expressions
nested in it.  It knows nothing of \NS, and it never fails on the output of
the lexer.  \emph{Resolution} then hands the pieces of each bracket to the
host parser, innermost markup expressions first, and yields the abstract
syntax.  All syntax errors other than lexical ones come from the host
parser.

\subsection{Raw nodes}

\begin{equation}\label{eq:raw-nodes}
\begin{aligned}
\rho &::= \kw{rblock}(u, s_u, \bar p, s) \mid \kw{rfinal}(u, s)
  && \text{raw node} \\
p &::= \kw{src}(w, s) \mid \kw{hole}(\bar\rho, s)
  && \text{piece}
\end{aligned}
\end{equation}
Here $s$ ranges over spans.  A raw block $\kw{rblock}(u, s_u, \bar p, s)$ is
a $\kw{DocumentBlock}$: its contents $u$, with escapes resolved and possibly
empty, the span $s_u$ of the contents, the pieces $\bar p$ of its bracket,
and the span $s$ of the whole block, from the start of its contents to the
end of its $\Close$.  A raw final node $\kw{rfinal}(u, s)$ is the trailing
$\kw{Text}$ when it is not empty.  A piece is either code $\kw{src}(w, s)$,
or $\kw{hole}(\bar\rho, s)$, a markup expression whose contents parse to
$\bar\rho$ and whose span $s$ runs from its $\Open$ to its $\Close$.

\subsection{The structural parse}

The judgment $\bar t \Rightarrow_{\mathrm{m}} (\bar\rho, \bar t')$ reads a
$\kw{DocumentMarkup}$ from the located tokens $\bar t$, leaving $\bar t'$,
which is empty or begins with $\kw{closeMarkup}$.  The judgment
$\bar t \Rightarrow_{\mathrm{b}} (\bar p, s, \bar t')$ reads the pieces of a
bracket up to and including its $\kw{closeMessages}$, whose span is $s$.
Tokens are written with their spans, as $\kw{text}(u)_{[i,j)}$.
\begin{gather}
\label{eq:markup-end}
\rulelabel{Markup-End}\quad
\epsilon \Rightarrow_{\mathrm{m}} (\epsilon, \epsilon)
\\
\label{eq:markup-stop}
\rulelabel{Markup-Stop}\quad
\kw{closeMarkup}_{s}\,\bar t \Rightarrow_{\mathrm{m}}
  (\epsilon, \kw{closeMarkup}_{s}\,\bar t)
\\
\label{eq:markup-trailing-end}
\rulelabel{Markup-Trailing-End}\quad
\kw{text}(u)_{s} \Rightarrow_{\mathrm{m}} (\kw{rfinal}(u, s), \epsilon)
\\
\label{eq:markup-trailing-stop}
\rulelabel{Markup-Trailing-Stop}\quad
\kw{text}(u)_{s}\,\kw{closeMarkup}_{s'}\,\bar t \Rightarrow_{\mathrm{m}}
  (\kw{rfinal}(u, s), \kw{closeMarkup}_{s'}\,\bar t)
\\
\label{eq:markup-block}
\rulelabel{Markup-Block}\quad
\frac{\bar t \Rightarrow_{\mathrm{b}} (\bar p, [k', l), \bar t')
      \qquad \bar t' \Rightarrow_{\mathrm{m}} (\bar\rho, \bar t'')}
     {\kw{text}(u)_{[i,j)}\,\kw{openMessages}_{[j,k)}\,\bar t
      \Rightarrow_{\mathrm{m}}
      (\kw{rblock}(u, [i,j), \bar p, [i,l))\,\bar\rho, \bar t'')}
\\
\label{eq:markup-bare-block}
\rulelabel{Markup-Bare-Block}\quad
\frac{\bar t \Rightarrow_{\mathrm{b}} (\bar p, [k', l), \bar t')
      \qquad \bar t' \Rightarrow_{\mathrm{m}} (\bar\rho, \bar t'')}
     {\kw{openMessages}_{[j,k)}\,\bar t
      \Rightarrow_{\mathrm{m}}
      (\kw{rblock}(\epsilon, [j,j), \bar p, [j,l))\,\bar\rho, \bar t'')}
\\
\label{eq:messages-close}
\rulelabel{Messages-Close}\quad
\kw{closeMessages}_{s}\,\bar t \Rightarrow_{\mathrm{b}} (\epsilon, s, \bar t)
\\
\label{eq:messages-code}
\rulelabel{Messages-Code}\quad
\frac{\bar t \Rightarrow_{\mathrm{b}} (\bar p, s', \bar t')}
     {\kw{code}(w)_{s}\,\bar t \Rightarrow_{\mathrm{b}}
      (\kw{src}(w, s)\,\bar p, s', \bar t')}
\\
\label{eq:messages-markup}
\rulelabel{Messages-Markup}\quad
\frac{\bar t \Rightarrow_{\mathrm{m}}
        (\bar\rho, \kw{closeMarkup}_{[k,l)}\,\bar t')
      \qquad \bar t' \Rightarrow_{\mathrm{b}} (\bar p, s', \bar t'')}
     {\kw{openMarkup}_{[i,j)}\,\bar t \Rightarrow_{\mathrm{b}}
      (\kw{hole}(\bar\rho, [i,l))\,\bar p, s', \bar t'')}
\end{gather}
A bracket with no text before it has empty contents, with the empty span at
the start of its $\Open$.  The trailing $\kw{Text}$ yields a node only when
it is not empty: an empty $\kw{Text}$ is the absence of a $\kw{text}$
token.  Because $\kw{text}$ tokens are never adjacent, the token after a
$\kw{text}$ token decides which rule applies, as the design notes argue.

\paragraph{Totality (obligation).}
\begin{equation}\label{eq:structural-totality}
W \text{ lexes to } \bar\ell
\implies
\exists!\, \bar\rho.\ \mathit{tokens}(\bar\ell) \Rightarrow_{\mathrm{m}}
  (\bar\rho, \epsilon)
\end{equation}
The lexer emits $\kw{text}$ and $\kw{openMessages}$ only in text mode,
$\kw{code}$ and $\kw{openMarkup}$ only in code mode, and its mode stack
pairs every $\kw{openMessages}$ with a $\kw{closeMessages}$ and every
$\kw{openMarkup}$ with a $\kw{closeMarkup}$, properly nested; no
$\kw{closeMarkup}$ occurs at top level, since that is
$\kw{strayClose}$.  So the structural parse cannot fail on lexer output.

\subsection{Abstract syntax}

\begin{equation}\label{eq:abstract-syntax}
\begin{aligned}
n &::= \kw{block}(u, s_u, \bar\sigma, s) \mid \kw{final}(u, s)
  && \text{DocumentBlockAST} \\
L &::= \kw{markup}(d, \bar n, s) && \text{MarkupLiteral} \\
\sigma &::= \kw{send}(\bar c, x, \bar\eta, s) \mid \kw{comment}(s)
  && \text{statement} \\
h &::= \kw{src}(w, s) \mid \kw{hole}(L) && \text{host input piece}
\end{aligned}
\end{equation}
A $\kw{block}$ is the DocumentBlockAST of a bracket: its contents and their
span, its statements in order, and its span.  A $\kw{final}$ node is the
final DocumentBlockAST of the design notes, the one with no messages; it is
kept distinct from a block whose bracket is empty, which is printed with
its bracket.  A markup literal holds the name $d$ of the source it was read
from, the DocumentBlockASTs of a markup expression, and its span.  Every
source has a name; a document's is fixed by the image it belongs to
(Section~\ref{sec:image}), and the name is what lets an edit find the source
of a literal that a method returned.  A statement is either a message send
$\kw{send}(\bar c, x, \bar\eta, s)$, with the spans $\bar c$ of the comments
that precede it in its bracket, its selector $x$, its argument expressions
$\bar\eta$, and its span; or a $\kw{comment}(s)$ that no message follows.
Argument expressions $\eta$ are abstract syntax of the host.  The core
observes only $\mathit{lits}(\eta)$, the markup literals occurring in
$\eta$, in source order.

\subsection{The host parser}
\label{sec:host-parser}

The host parser is a function from the host input pieces of one bracket to
either its statements or a syntax error with its span:
\begin{equation}\label{eq:host-parser}
\mathcal{P}(\bar h) \in \{\, \kw{stmts}(\bar\sigma) \,\} \cup
  \{\, \kw{fail}(s) \,\}
\end{equation}
A hole is an atomic expression, a literal, to the host parser.  The host
parser is required to satisfy:
\begin{enumerate}
\item[(H1)] it is total: parsing does not diverge;
\item[(H2)] every send is a message without an explicit receiver, unary or
      keyword, and not a cascade; any other statement makes it fail;
\item[(H3)] the holes of $\bar h$, in order, are exactly the concatenation
      of $\mathit{lits}(\eta)$ over the arguments of the sends, in order, so
      no markup expression is lost, duplicated, or reordered;
\item[(H4)] statement, comment, and error spans lie within the span from
      the first piece of $\bar h$ to its last, statements in increasing
      order; and
\item[(H5)] code that is empty or white space has no statements; a
      bracket holding only a comment yields one $\kw{comment}$ statement;
      a comment followed by a message in its bracket is recorded in that
      send's $\bar c$.
\end{enumerate}
Ordinary \NS{} supplies $\mathcal{P}$ from its front-end semantics; the
obligation that $\Open$ and $\Close$ never arise in valid \NS{} code
(Section~\ref{sec:lexical-errors}) is what lets the core cut the code at
the holes.

\subsection{Resolution}

Resolution replaces each raw node by a DocumentBlockAST, innermost markup
expressions first, and stops at the first syntax error.  A result is
$\kw{ok}(x)$ or $\kw{fail}(s)$.  As for the readers, write
$v \cdot \kw{ok}(\bar x) = \kw{ok}(v\,\bar x)$ and
$v \cdot \kw{fail}(s) = \kw{fail}(s)$.
\begin{equation}\label{eq:resolve}
\begin{aligned}
\mathit{res}(\epsilon) &= \kw{ok}(\epsilon) \\
\mathit{res}(\rho\,\bar\rho) &= n \cdot \mathit{res}(\bar\rho)
  && \text{if } \mathit{node}(\rho) = \kw{ok}(n) \\
\mathit{res}(\rho\,\bar\rho) &= \kw{fail}(s)
  && \text{if } \mathit{node}(\rho) = \kw{fail}(s) \\[1ex]
\mathit{node}(\kw{rfinal}(u, s)) &= \kw{ok}(\kw{final}(u, s)) \\
\mathit{node}(\kw{rblock}(u, s_u, \bar p, s)) &=
  \kw{ok}(\kw{block}(u, s_u, \bar\sigma, s))
  && \text{if } \mathit{pieces}(\bar p) = \kw{ok}(\bar h),\
     \mathcal{P}(\bar h) = \kw{stmts}(\bar\sigma) \\
\mathit{node}(\kw{rblock}(u, s_u, \bar p, s)) &= \kw{fail}(s')
  && \text{if } \mathit{pieces}(\bar p) = \kw{ok}(\bar h),\
     \mathcal{P}(\bar h) = \kw{fail}(s') \\
\mathit{node}(\kw{rblock}(u, s_u, \bar p, s)) &= \kw{fail}(s')
  && \text{if } \mathit{pieces}(\bar p) = \kw{fail}(s') \\[1ex]
\mathit{pieces}(\epsilon) &= \kw{ok}(\epsilon) \\
\mathit{pieces}(\kw{src}(w, s)\,\bar p) &=
  \kw{src}(w, s) \cdot \mathit{pieces}(\bar p) \\
\mathit{pieces}(\kw{hole}(\bar\rho, s)\,\bar p) &=
  \kw{hole}(L) \cdot \mathit{pieces}(\bar p)
  && \text{if } \mathit{res}(\bar\rho) = \kw{ok}(\bar n),\
     L = \kw{markup}(d, \bar n, s) \\
\mathit{pieces}(\kw{hole}(\bar\rho, s)\,\bar p) &= \kw{fail}(s')
  && \text{if } \mathit{res}(\bar\rho) = \kw{fail}(s')
\end{aligned}
\end{equation}
Here $d$ is the name of the source being parsed, which every literal records.
The error reported is the first in this order: within a bracket, an error
in a nested markup expression is found before the host parses the bracket
itself.  Every markup expression is therefore parsed when the enclosing
markup is parsed, as the design notes require; only its interpretation is
deferred (Section~\ref{sec:interpretation}).

\subsection{Parsing a source}

\begin{equation}\label{eq:parse-results}
\pi ::= \kw{parsed}(\bar n) \mid \kw{lexError}(e, i)
  \mid \kw{syntaxError}(s)
\end{equation}
\begin{gather}
\label{eq:parse-lex-error}
\rulelabel{Parse-Lex-Error}\quad
\frac{W \text{ fails with } e \text{ at } i}
     {\mathit{parse}(W) = \kw{lexError}(e, i)}
\\
\label{eq:parse-ok}
\rulelabel{Parse-Ok}\quad
\frac{W \text{ lexes to } \bar\ell
      \qquad \mathit{tokens}(\bar\ell) \Rightarrow_{\mathrm{m}}
        (\bar\rho, \epsilon)
      \qquad \mathit{res}(\bar\rho) = \kw{ok}(\bar n)}
     {\mathit{parse}(W) = \kw{parsed}(\bar n)}
\\
\label{eq:parse-syntax-error}
\rulelabel{Parse-Syntax-Error}\quad
\frac{W \text{ lexes to } \bar\ell
      \qquad \mathit{tokens}(\bar\ell) \Rightarrow_{\mathrm{m}}
        (\bar\rho, \epsilon)
      \qquad \mathit{res}(\bar\rho) = \kw{fail}(s)}
     {\mathit{parse}(W) = \kw{syntaxError}(s)}
\end{gather}
By the determinism of lexing, structural totality, and (H1), every source
has exactly one parse result.

\paragraph{Provisional: one error fails the parse.}
A syntax error anywhere makes the whole source fail to parse.  That is
accepted for now (2026-11-08) but is not the intended end state: an editor
that reparses on every keystroke will want local recovery, in which a
bracket the host cannot parse becomes an error leaf and the rest of the
document still parses, as a message that is not understood does in
interpretation.

\paragraph{Example.}
The source $\str{[\bs{}title: [\bs{}Hamlet\bs]\bs]}$ lexes to the tokens
\begin{gather*}
\kw{openMessages}_{[0,2)}\ \kw{code}(\str{title: })_{[2,9)}\
\kw{openMarkup}_{[9,11)}
\\
\kw{text}(\str{Hamlet})_{[11,17)}\
\kw{closeMarkup}_{[17,19)}\ \kw{closeMessages}_{[19,21)}
\end{gather*}
whose structural parse is the single raw block
\[
\kw{rblock}(\epsilon, [0,0),\
  \kw{src}(\str{title: }, [2,9))\
  \kw{hole}(\kw{rfinal}(\str{Hamlet}, [11,17)), [9,19)),\ [0,21)).
\]
Resolution gives the hole the literal
$L = \kw{markup}(d, \kw{final}(\str{Hamlet}, [11,17)), [9,19))$, and the host
parser yields one send of $\str{title:}$ whose single argument $\eta$ has
$\mathit{lits}(\eta) = L$.

\section{Interpretation}
\label{sec:interpretation}

Interpretation turns a sequence of DocumentBlockASTs into a tree of
document blocks.  As in the design notes, it keeps a stack of open blocks,
initially holding only a root, and processes the ASTs in order: the
contents of each are added to the block on top of the stack (the ToS)
under the whitespace rules, and its messages are then sent, one at a time,
to whichever block is ToS when each is sent.  Each message pops the stack,
adds a leaf, adds a block and pushes it, mutates the ToS, or is not
understood.  Blocks are attached when they are created, so every block on
the stack is already part of the tree.

\subsection{The host evaluator and the block signature}

The host evaluator and the block signature are parameters.  Neither is
defined by the core; the core fixes only their shape.

\paragraph{Host evaluator.}
Let $\chi$ range over host states, $v$ over values, and $\xi$ over
exceptions.  Host states are abstract; they include the document, which
arguments are evaluated against as their implicit receiver.  The
evaluator relates a host state and an argument expression to an outcome:
\begin{equation}\label{eq:host-evaluator}
\langle \chi, \eta \rangle \Downarrow_{\mathrm{E}} o
\qquad
o ::= \kw{val}(\chi', v) \mid \kw{exn}(\chi', \xi)
\end{equation}
Evaluation may diverge, in which case there is no outcome.  The evaluator
is required to satisfy:
\begin{enumerate}
\item[(E1)] a markup expression standing alone as an argument evaluates,
      without effect, to the uninterpreted literal: if $\eta$ is the hole
      $L$ alone, then
      $\langle \chi, \eta \rangle \Downarrow_{\mathrm{E}}
        \kw{val}(\chi, \ulcorner L \urcorner)$; and
\item[(E2)] values never refer to the blocks the interpreter is building.
      A value may contain a finished tree $\tau$
      (Section~\ref{sec:interp-results}), which is how a computed block
      reaches the markup.  A tree held in a value is written $\tau^{L}$:
      it records the literal $L$ it was interpreted from, which is where an
      edit of the tree goes (Section~\ref{sec:editing}).
\end{enumerate}
Values thus include the literals $\ulcorner L \urcorner$ and the trees
$\tau^{L}$; everything else about them belongs to the host.

\paragraph{Block signature.}
Let $\kappa$ range over block states.  A block state is the host's view of
a block: its class and its fields, such as the title of the top-level
block.  The signature gives, for each block state, its \emph{content
class}
\begin{equation}\label{eq:content-class}
\mathit{cls}(\kappa) \in \{\, \kw{flow}, \kw{para}, \kw{cell}, \kw{inline} \,\}
\end{equation}
and, for a message sent to a block, its effect:
\begin{equation}\label{eq:block-signature}
\begin{gathered}
\langle \chi, \kappa, x, \bar v \rangle \Downarrow_{\mathrm{S}} o
\qquad
o ::= \kw{eff}(\chi', \varepsilon) \mid \kw{exn}(\chi', \xi)
\\
\varepsilon ::= \kw{pop} \mid \kw{popReject} \mid \kw{leaf}(\kappa')
  \mid \kw{push}(\kappa') \mid \kw{mutate}(\kappa') \mid \kw{refuse}
  \mid \kw{dnu}
\end{gathered}
\end{equation}
Here $x$ is the selector and $\bar v$ the argument values.  A $\kw{flow}$
block holds no inline content: the top-level block, a section, a list, a
table or row, a quote, a figure, a list item.  A $\kw{para}$ block holds
inline content and is block-level, so its text is trimmed when it closes: a
paragraph.  A $\kw{cell}$ block, a table cell, holds either flow content or
inline content, but not both; which one is decided when it closes.  An
$\kw{inline}$ block holds inline content and keeps its white space: a
style, or an inline root.  A message that the block does not understand has
effect $\kw{dnu}$.  A message that the block understands but whose check
fails, such as a second title or an argument of the wrong kind, has effect
$\kw{refuse}$: it leaves an error leaf and changes nothing else.  A closing
message whose check fails has effect $\kw{popReject}$: the block is popped
and closed as by $\kw{pop}$, and flagged as rejected.  An ordinary error
raised in the method that implements a message, including a failed check
that the method chooses to raise rather than report, is an exception.

The effect relation may itself interpret markup: a method such as
$\str{title:}$ receives a literal $\ulcorner L \urcorner$ and interprets it
with a root of its choosing (Section~\ref{sec:interp-results}).  The
signature and interpretation are therefore defined together, as the least
relations closed under the rules of the signature and of this section.

A block state is the only part of a block the host sees.  The children of
a block, and the messages it received, are kept by the core in the block
store.

\subsection{White space}

White space is the set
$\mathit{WS} = \{\text{U+0009}, \text{U+000A}, \text{U+000B},
\text{U+000C}, \text{U+000D}, \text{U+0020}\}$; in particular a
non-breaking space is not white space, so it is kept as written
(whitespace rule~1).  Write $\vs$ for U+0020.
$\mathit{collapse}(u)$ replaces each maximal run of characters of
$\mathit{WS}$ in $u$ by a single $\vs$:
\begin{equation}\label{eq:collapse}
\begin{aligned}
\mathit{collapse}(\epsilon) &= \epsilon \\
\mathit{collapse}(c\,u) &=
  \vs\ \mathit{collapse}(\mathit{dropWS}(u))
  && \text{if } c \in \mathit{WS} \\
\mathit{collapse}(c\,u) &= c\ \mathit{collapse}(u)
  && \text{if } c \notin \mathit{WS}
\end{aligned}
\end{equation}
where $\mathit{dropWS}(u)$ removes the leading characters of $u$ that are in
$\mathit{WS}$.

The contents $u$, with span $s$, of a DocumentBlockAST add at most one
child to the ToS, decided by the ToS's content class $k$
(whitespace rules~3 and~4):
\begin{equation}\label{eq:text-child}
\mathit{child}(k, u, s) =
\begin{cases}
\kw{none} & \text{if } \mathit{collapse}(u) = \epsilon \\
\kw{none} & \text{if } \mathit{collapse}(u) = \vs
  \text{ and } k = \kw{flow} \\
\kw{run}(\vs, s) & \text{if } \mathit{collapse}(u) =
  \vs \text{ and } k \neq \kw{flow} \\
\kw{stray}(u, s) & \text{if } \mathit{collapse}(u) \notin
  \{\epsilon, \vs\} \text{ and } k = \kw{flow} \\
\kw{run}(\mathit{collapse}(u), s) & \text{if } \mathit{collapse}(u)
  \notin \{\epsilon, \vs\} \text{ and } k \neq \kw{flow}
\end{cases}
\end{equation}
A text run keeps the span of the contents it came from.  A $\kw{stray}$
child is the error leaf of rule~4; it keeps the contents as written.  A
$\kw{cell}$ block keeps its text as a $\kw{para}$ block does; whether the
text may stay is decided when the cell closes.

When a $\kw{para}$ block closes, its text is trimmed (whitespace rule~2):
a leading $\vs$ is removed from its first child and a
trailing one from its last child, when those children are text runs, and
a run left empty is removed.
\begin{equation}\label{eq:trim}
\begin{aligned}
\mathit{trim}(\bar g) &= \mathit{trimEnd}(\mathit{trimStart}(\bar g)) \\
\mathit{trimStart}(\kw{run}(\vs, s)\,\bar g) &= \bar g \\
\mathit{trimStart}(\kw{run}(\vs\,c\,u, s)\,\bar g) &=
  \kw{run}(c\,u, s)\,\bar g \\
\mathit{trimStart}(\bar g) &= \bar g
  && \text{if } \bar g \neq \kw{run}(\vs\,u, s)\,\bar g' \\
\mathit{trimEnd}(\bar g\,\kw{run}(\vs, s)) &= \bar g \\
\mathit{trimEnd}(\bar g\,\kw{run}(u\,c\,\vs, s)) &=
  \bar g\,\kw{run}(u\,c, s) \\
\mathit{trimEnd}(\bar g) &= \bar g
  && \text{if } \bar g \neq \bar g'\,\kw{run}(u\,\vs, s)
\end{aligned}
\end{equation}
Since runs are collapsed, one space is all there is to trim at each end.
White space at the edge of an inline child, such as an emphasized run that
begins with a space, is the inline block's own and is kept.

\subsection{Block store and configurations}

\begin{equation}\label{eq:block-store}
\begin{aligned}
g &::= \kw{run}(u, s) \mid \kw{blk}(\beta) \mid \kw{note}(s)
  \mid \kw{mut}(\sigma) \mid \kw{dnu}(\sigma) \\
  &\phantom{{}::={}} \mid \kw{stray}(u, s)
  \mid \kw{rootPop}(\sigma) \mid \kw{refused}(\sigma)
  && \text{child} \\
B &::= \langle \kappa, \bar g, \bar\sigma, r \rangle
  \qquad r ::= \kw{ok} \mid \kw{rejected}
  && \text{block record} \\
\Sigma &::= B_0 \cdots B_{m-1}
  && \text{block store} \\
\omega &::= \kw{text}(u, s) \mid \sigma && \text{item} \\
C &::= \langle \chi, \Sigma, \bar\beta, \bar\omega \rangle
  \mid \kw{raise}(\chi, \xi)
  && \text{configuration}
\end{aligned}
\end{equation}
A block reference $\beta$ is an index into the store, and a new block takes
the next index, $|\Sigma|$.  A block record holds the block's state, its
children in order, and the messages it received, in order: the one that
created it, any that mutated it, and the one that closed it, and whether
it was rejected when it closed.  Children are
text runs, blocks, the $\kw{note}$ of a comment bracket (a CommentBlock),
the marker $\kw{mut}(\sigma)$ of a mutating message, which records where
among the children the block received it and renders nothing, and four
error leaves: an ErrorDocumentBlock $\kw{dnu}(\sigma)$ holding the
message not understood, stray text, $\kw{rootPop}(\sigma)$, a message
that would pop the root, and $\kw{refused}(\sigma)$, a message whose check
failed.  The stack $\bar\beta$ is written top first and is
never empty; its last element is the root.

The items of a DocumentBlockAST are its contents followed by its
statements:
\begin{equation}\label{eq:items}
\mathit{items}(\kw{block}(u, s_u, \bar\sigma, s)) = \kw{text}(u, s_u)\,\bar\sigma
\qquad
\mathit{items}(\kw{final}(u, s)) = \kw{text}(u, s)
\end{equation}
and $\mathit{items}(\bar n)$ concatenates the items of $\bar n$.

On the store, $\Sigma[\beta \mathrel{+}= g]$ appends the child $g$ to block
$\beta$, $\Sigma[\beta \mathrel{\lhd} \sigma]$ appends the message $\sigma$
to its received messages, $\Sigma[\beta := \kappa]$ replaces its state,
$\Sigma[\beta\ \kw{rejected}]$ flags it as rejected, and
$\Sigma\,B$ allocates $B$ at index $|\Sigma|$.

In a store $\Sigma$, a child $\kw{blk}(\beta')$ is \emph{inline} if
$\mathit{cls}(\kappa') = \kw{inline}$ and a \emph{block child} otherwise,
where $\kappa'$ is the state of $\Sigma(\beta')$.  Closing a block trims
it when it is a $\kw{para}$ block, or a $\kw{cell}$ block with no block
child.  A cell with a block child holds flow content: its white-space runs
are layout and are dropped, and if text or an inline child remains, the
cell mixes the two and is rejected.  Let
$\Sigma(\beta) = \langle \kappa, \bar g, \bar\sigma, r\rangle$.
\begin{equation}\label{eq:close}
\mathit{close}(\Sigma, \beta) =
\begin{cases}
\Sigma[\beta \mapsto \langle \kappa, \mathit{trim}(\bar g), \bar\sigma, r\rangle]
  & \text{if } \mathit{cls}(\kappa) = \kw{para}, \text{ or }
    \mathit{cls}(\kappa) = \kw{cell} \\
  & \quad\text{and } \bar g \text{ has no block child} \\
\Sigma[\beta \mapsto \langle \kappa, \bar g', \bar\sigma, r'\rangle]
  & \text{if } \mathit{cls}(\kappa) = \kw{cell}
    \text{ and } \bar g \text{ has a block child} \\
\Sigma & \text{otherwise}
\end{cases}
\end{equation}
where $\bar g'$ is $\bar g$ without its runs $\kw{run}(\vs, s)$, and
$r' = \kw{rejected}$ if $\bar g'$ has a run or an inline child, and
$r' = r$ otherwise.

\subsection{Argument evaluation}

Arguments are evaluated left to right, threading the host state; the first
exception ends the evaluation.
\begin{gather}
\label{eq:args-nil}
\rulelabel{Args-Nil}\quad
\langle \chi, \epsilon \rangle \Downarrow_{\mathrm{A}} \kw{vals}(\chi, \epsilon)
\\
\label{eq:args-cons}
\rulelabel{Args-Cons}\quad
\frac{\langle \chi, \eta \rangle \Downarrow_{\mathrm{E}} \kw{val}(\chi', v)
      \qquad \langle \chi', \bar\eta \rangle \Downarrow_{\mathrm{A}}
        \kw{vals}(\chi'', \bar v)}
     {\langle \chi, \eta\,\bar\eta \rangle \Downarrow_{\mathrm{A}}
        \kw{vals}(\chi'', v\,\bar v)}
\\
\label{eq:args-raise-head}
\rulelabel{Args-Raise-Head}\quad
\frac{\langle \chi, \eta \rangle \Downarrow_{\mathrm{E}} \kw{exn}(\chi', \xi)}
     {\langle \chi, \eta\,\bar\eta \rangle \Downarrow_{\mathrm{A}}
        \kw{exn}(\chi', \xi)}
\\
\label{eq:args-raise-tail}
\rulelabel{Args-Raise-Tail}\quad
\frac{\langle \chi, \eta \rangle \Downarrow_{\mathrm{E}} \kw{val}(\chi', v)
      \qquad \langle \chi', \bar\eta \rangle \Downarrow_{\mathrm{A}}
        \kw{exn}(\chi'', \xi)}
     {\langle \chi, \eta\,\bar\eta \rangle \Downarrow_{\mathrm{A}}
        \kw{exn}(\chi'', \xi)}
\end{gather}
A message not understood while computing an argument is sent to some
object other than the ToS, so it surfaces here as an exception, an ordinary
error rather than a markup error, as the design notes require.

\subsection{Steps}

The step relation $C \longrightarrow C'$ processes one item.  Throughout,
$\Sigma(\beta) = \langle \kappa, \bar g, \bar\sigma', r \rangle$ is the
record of the ToS.  In the rules for a send,
$\sigma = \kw{send}(\bar c, x, \bar\eta, s)$, and the premises
\[
\langle \chi, \bar\eta \rangle \Downarrow_{\mathrm{A}} \kw{vals}(\chi', \bar v)
\qquad
\langle \chi', \kappa, x, \bar v \rangle \Downarrow_{\mathrm{S}}
  \kw{eff}(\chi'', \varepsilon)
\]
are abbreviated as $\chi, \beta \vdash \sigma \Downarrow \chi'', \varepsilon$.
\begin{gather}
\label{eq:step-text}
\rulelabel{Step-Text}\quad
\frac{\mathit{child}(\mathit{cls}(\kappa), u, s) = g}
     {\langle \chi, \Sigma, \beta\,\bar\beta, \kw{text}(u, s)\,\bar\omega \rangle
      \longrightarrow
      \langle \chi, \Sigma[\beta \mathrel{+}= g], \beta\,\bar\beta, \bar\omega \rangle}
\\
\label{eq:step-text-none}
\rulelabel{Step-Text-None}\quad
\frac{\mathit{child}(\mathit{cls}(\kappa), u, s) = \kw{none}}
     {\langle \chi, \Sigma, \beta\,\bar\beta, \kw{text}(u, s)\,\bar\omega \rangle
      \longrightarrow
      \langle \chi, \Sigma, \beta\,\bar\beta, \bar\omega \rangle}
\\
\label{eq:step-comment}
\rulelabel{Step-Comment}\quad
\langle \chi, \Sigma, \beta\,\bar\beta, \kw{comment}(s)\,\bar\omega \rangle
\longrightarrow
\langle \chi, \Sigma[\beta \mathrel{+}= \kw{note}(s)], \beta\,\bar\beta,
  \bar\omega \rangle
\\
\label{eq:step-pop}
\rulelabel{Step-Pop}\quad
\frac{\chi, \beta \vdash \sigma \Downarrow \chi'', \kw{pop}}
     {\langle \chi, \Sigma, \beta\,\beta'\,\bar\beta, \sigma\,\bar\omega \rangle
      \longrightarrow
      \langle \chi'', \mathit{close}(\Sigma[\beta \lhd \sigma], \beta),
        \beta'\,\bar\beta, \bar\omega \rangle}
\\
\label{eq:step-pop-reject}
\rulelabel{Step-Pop-Reject}\quad
\frac{\chi, \beta \vdash \sigma \Downarrow \chi'', \kw{popReject}}
     {\langle \chi, \Sigma, \beta\,\beta'\,\bar\beta, \sigma\,\bar\omega \rangle
      \longrightarrow
      \langle \chi'', \mathit{close}(\Sigma[\beta \lhd \sigma], \beta)
        [\beta\ \kw{rejected}], \beta'\,\bar\beta, \bar\omega \rangle}
\\
\label{eq:step-root-pop}
\rulelabel{Step-Root-Pop}\quad
\frac{\chi, \beta \vdash \sigma \Downarrow \chi'', \varepsilon
      \qquad \varepsilon \in \{\kw{pop}, \kw{popReject}\}}
     {\langle \chi, \Sigma, \beta, \sigma\,\bar\omega \rangle
      \longrightarrow
      \langle \chi'', \Sigma[\beta \mathrel{+}= \kw{rootPop}(\sigma)], \beta,
        \bar\omega \rangle}
\\
\label{eq:step-leaf}
\rulelabel{Step-Leaf}\quad
\frac{\chi, \beta \vdash \sigma \Downarrow \chi'', \kw{leaf}(\kappa')
      \qquad \beta'' = |\Sigma|}
     {\langle \chi, \Sigma, \beta\,\bar\beta, \sigma\,\bar\omega \rangle
      \longrightarrow
      \langle \chi'', \Sigma[\beta \mathrel{+}= \kw{blk}(\beta'')]\,
        \langle \kappa', \epsilon, \sigma, \kw{ok} \rangle,
        \beta\,\bar\beta, \bar\omega \rangle}
\\
\label{eq:step-push}
\rulelabel{Step-Push}\quad
\frac{\chi, \beta \vdash \sigma \Downarrow \chi'', \kw{push}(\kappa')
      \qquad \beta'' = |\Sigma|}
     {\langle \chi, \Sigma, \beta\,\bar\beta, \sigma\,\bar\omega \rangle
      \longrightarrow
      \langle \chi'', \Sigma[\beta \mathrel{+}= \kw{blk}(\beta'')]\,
        \langle \kappa', \epsilon, \sigma, \kw{ok} \rangle,
        \beta''\,\beta\,\bar\beta, \bar\omega \rangle}
\\
\label{eq:step-mutate}
\rulelabel{Step-Mutate}\quad
\frac{\chi, \beta \vdash \sigma \Downarrow \chi'', \kw{mutate}(\kappa')}
     {\langle \chi, \Sigma, \beta\,\bar\beta, \sigma\,\bar\omega \rangle
      \longrightarrow
      \langle \chi'', \Sigma[\beta := \kappa'][\beta \lhd \sigma]
        [\beta \mathrel{+}= \kw{mut}(\sigma)],
        \beta\,\bar\beta, \bar\omega \rangle}
\\
\label{eq:step-dnu}
\rulelabel{Step-Dnu}\quad
\frac{\chi, \beta \vdash \sigma \Downarrow \chi'', \kw{dnu}}
     {\langle \chi, \Sigma, \beta\,\bar\beta, \sigma\,\bar\omega \rangle
      \longrightarrow
      \langle \chi'', \Sigma[\beta \mathrel{+}= \kw{dnu}(\sigma)],
        \beta\,\bar\beta, \bar\omega \rangle}
\\
\label{eq:step-refuse}
\rulelabel{Step-Refuse}\quad
\frac{\chi, \beta \vdash \sigma \Downarrow \chi'', \kw{refuse}}
     {\langle \chi, \Sigma, \beta\,\bar\beta, \sigma\,\bar\omega \rangle
      \longrightarrow
      \langle \chi'', \Sigma[\beta \mathrel{+}= \kw{refused}(\sigma)],
        \beta\,\bar\beta, \bar\omega \rangle}
\\
\label{eq:step-arg-raise}
\rulelabel{Step-Arg-Raise}\quad
\frac{\langle \chi, \bar\eta \rangle \Downarrow_{\mathrm{A}} \kw{exn}(\chi', \xi)}
     {\langle \chi, \Sigma, \beta\,\bar\beta, \sigma\,\bar\omega \rangle
      \longrightarrow \kw{raise}(\chi', \xi)}
\\
\label{eq:step-effect-raise}
\rulelabel{Step-Effect-Raise}\quad
\frac{\langle \chi, \bar\eta \rangle \Downarrow_{\mathrm{A}} \kw{vals}(\chi', \bar v)
      \qquad \langle \chi', \kappa, x, \bar v \rangle \Downarrow_{\mathrm{S}}
        \kw{exn}(\chi'', \xi)}
     {\langle \chi, \Sigma, \beta\,\bar\beta, \sigma\,\bar\omega \rangle
      \longrightarrow \kw{raise}(\chi'', \xi)}
\end{gather}
In Step-Pop and Step-Pop-Reject the stack has at least two blocks; in
Step-Root-Pop it has only the root.  The root is never popped: a message that would pop it leaves an
error leaf and changes nothing else, so, as with a message not understood,
later messages behave as though it were absent.  No message both pops and
pushes, since an effect is exactly one of the seven.

\subsection{Results}
\label{sec:interp-results}

A finished tree is
\begin{equation}\label{eq:trees}
\begin{aligned}
\tau &::= \kw{tree}(\kappa, \bar t, \bar\sigma, f)
  && f ::= \kw{closed} \mid \kw{unclosed} \mid \kw{rejected} \\
t &::= \kw{run}(u, s) \mid \kw{sub}(\tau) \mid \kw{note}(s)
  \mid \kw{mut}(\sigma) \mid \kw{dnu}(\sigma) \\
  &\phantom{{}::={}} \mid \kw{stray}(u, s)
  \mid \kw{rootPop}(\sigma) \mid \kw{refused}(\sigma)
\end{aligned}
\end{equation}
and is read off the store from the root, flagging the blocks in a set $U$
as unclosed:
\begin{equation}\label{eq:read-tree}
\mathit{tree}_U(\Sigma, \beta) =
  \kw{tree}(\kappa, \mathit{tree}_U(\Sigma, \bar g), \bar\sigma, f)
\quad\text{where } \Sigma(\beta) = \langle \kappa, \bar g, \bar\sigma, r \rangle
\end{equation}
and $f = \kw{unclosed}$ if $\beta \in U$, $f = \kw{rejected}$ if
$\beta \notin U$ and $r = \kw{rejected}$, and $f = \kw{closed}$ otherwise,
mapping $\kw{blk}(\beta')$ to $\kw{sub}(\mathit{tree}_U(\Sigma, \beta'))$
and every other child to itself.  Every block other than the root is
allocated as a child of a block allocated before it, so the store is a
tree rooted at index 0 and $\mathit{tree}$ is well defined.

Interpretation from a root state $\kappa_0$ runs the configuration
$\langle \chi, \langle \kappa_0, \epsilon, \epsilon, \kw{ok} \rangle, 0,
\mathit{items}(\bar n) \rangle$ to the end of its items:
\begin{equation}\label{eq:interp-results}
\psi ::= \kw{done}(\chi, \tau) \mid \kw{raised}(\chi, \xi)
\end{equation}
\begin{gather}
\label{eq:interp-done}
\rulelabel{Interp-Done}\quad
\langle \chi, \Sigma, \bar\beta\,0, \epsilon \rangle \Downarrow
  \kw{done}(\chi, \mathit{tree}_{\{\bar\beta\}}(\Sigma, 0))
\\
\label{eq:interp-step}
\rulelabel{Interp-Step}\quad
\frac{C \longrightarrow \langle \chi', \Sigma', \bar\beta', \bar\omega' \rangle
      \qquad \langle \chi', \Sigma', \bar\beta', \bar\omega' \rangle \Downarrow \psi}
     {C \Downarrow \psi}
\\
\label{eq:interp-raise}
\rulelabel{Interp-Raise}\quad
\frac{C \longrightarrow \kw{raise}(\chi', \xi)}
     {C \Downarrow \kw{raised}(\chi', \xi)}
\\
\label{eq:interpret}
\rulelabel{Interpret}\quad
\frac{\langle \chi, \langle \kappa_0, \epsilon, \epsilon, \kw{ok} \rangle, 0,
        \mathit{items}(\bar n) \rangle \Downarrow \psi}
     {\langle \chi, \kappa_0, \bar n \rangle \Downarrow_{\mathrm{I}} \psi}
\end{gather}
In Interp-Done the blocks still on the stack above the root were never
closed.  They are already in the tree, and each is flagged
$\kw{unclosed}$, so it can render plausibly from its subtrees; it is not
trimmed.  A literal is interpreted with a root chosen by the method that
receives it; this is what the design notes write as
\texttt{markup interpretWith: root}:
\begin{equation}\label{eq:interpret-with}
\mathit{interpretWith}(\chi, \kw{markup}(d, \bar n, s), \kappa_0) = \psi
\quad\text{iff}\quad
\langle \chi, \kappa_0, \bar n \rangle \Downarrow_{\mathrm{I}} \psi
\end{equation}
A document whose source parses to $\bar n$ is interpreted with the
signature's top-level block state $\kappa_{\mathrm{doc}}$, the
document's ToplevelDocumentBlock, as its root.

\paragraph{Errors.}
The markup errors of a tree are its error leaves and its unclosed and
rejected blocks, in document order:
\begin{equation}\label{eq:errors}
\begin{aligned}
\mathit{errors}(\kw{tree}(\kappa, \bar t, \bar\sigma, f)) &=
  \mathit{flag}(f)\ \mathit{errors}(t_1) \cdots \mathit{errors}(t_m) \\
\mathit{flag}(\kw{unclosed}) &= \kw{unclosed}(\bar\sigma),\quad
\mathit{flag}(\kw{rejected}) = \kw{rejected}(\bar\sigma),\quad
\mathit{flag}(\kw{closed}) = \epsilon \\
\mathit{errors}(\kw{sub}(\tau)) &= \mathit{errors}(\tau) \\
\mathit{errors}(t) &= t
  \quad\text{if } t \text{ is a } \kw{dnu},\ \kw{stray},\ \kw{rootPop},
    \text{ or } \kw{refused} \text{ leaf} \\
\mathit{errors}(t) &= \epsilon
  \quad\text{if } t \text{ is a } \kw{run},\ \kw{note},
    \text{ or } \kw{mut}
\end{aligned}
\end{equation}
An unclosed or rejected block is reported by the messages it received, the
first of which created it.  Markup errors never stop interpretation; only an
exception from the host does.  Errors inside a tree built by a nested
interpretation are the business of the method that stored it; the
standard library reports them with the tree that holds them.

\paragraph{Properties.}
Each step consumes one item, so a derivation of
$\langle \chi, \kappa_0, \bar n \rangle \Downarrow_{\mathrm{I}} \psi$ has
at most $|\mathit{items}(\bar n)|$ steps, not counting those of nested
interpretations, and interpretation can fail to
have a result only if the host evaluator or the signature diverges.
\begin{multline}\label{eq:interp-determinism}
\text{if } \Downarrow_{\mathrm{E}} \text{ and } \Downarrow_{\mathrm{S}}
\text{ each relate an input to at most one outcome,} \\
\text{then }
\langle \chi, \kappa_0, \bar n \rangle \Downarrow_{\mathrm{I}} \psi
\text{ for at most one } \psi
\end{multline}
The side conditions of the step rules are disjoint: an item is text, a
comment, or a send; the effect of a send is one of five; and Step-Pop and
Step-Root-Pop are told apart by the height of the stack.

\paragraph{Example.}
Take a signature in which $\str{paragraph}$, sent to a $\kw{flow}$ block,
pushes a $\kw{para}$ block, $\str{emphasis}$ sent to a $\kw{para}$ or
$\kw{inline}$ block pushes an $\kw{inline}$ one, and $\str{endParagraph}$
and $\str{endEmphasis}$ pop.  The source
\[
\str{[\bs{}paragraph\bs]\ Hi\ [\bs{}emphasis\bs]all[\bs{}endEmphasis\bs]\ [\bs{}endParagraph\bs]}
\]
has four DocumentBlockASTs, and the root $\beta = 0$ is a $\kw{flow}$
block.  The first has empty contents, which add nothing, and pushes a
paragraph, block~1.  The second adds the run $\vs\str{Hi}
\vs$ to block~1 and pushes the emphasis, block~2.  The third
adds the run $\str{all}$ to block~2 and pops it.  The fourth adds the
whitespace-only run $\vs$ to block~1, since a paragraph keeps
it, and pops the paragraph, which trims it: the leading space of the first
run goes, and the last run, now empty, is removed.  The tree is
\[
\kw{tree}(\kappa_{\mathrm{doc}},\ 
  \kw{sub}(\kw{tree}(\kappa_{\mathrm{p}},\
    \kw{run}(\str{Hi}\vs, s_2)\
    \kw{sub}(\kw{tree}(\kappa_{\mathrm{e}}, \kw{run}(\str{all}, s_3),
      \sigma_2\,\sigma_3, \kw{closed})),\
    \sigma_1\,\sigma_4, \kw{closed})),\
  \epsilon, \kw{closed})
\]
where $\sigma_1, \ldots, \sigma_4$ are the four sends and $s_2, s_3$ the
spans of the contents of the second and third ASTs.  The space between
\str{Hi} and the emphasis survives, as whitespace rule~2 intends.

\section{Printing}
\label{sec:printing}

The design notes ask that a document be a tree that prints itself to a
markup program reproducing it.  The printer works from a tree $\tau$ and the
source $W$ that it was interpreted from.  Text is printed from the runs, so
white space comes out normalized.  Messages are printed by the host's
pretty printer, a further parameter of the semantics, which for \NS{} is
the \NS{} pretty printer; the core prints only the markup literals in the
messages' arguments.  A
tree built by code has no source of its own; it reaches a document only as
a value, through the source of the message that computed it, and so needs
no printing of its own here.

\subsection{Escaping text}

\begin{equation}\label{eq:escape}
\begin{aligned}
\mathit{esc}(\epsilon) &= \epsilon \\
\mathit{esc}(\str{\bs}\,u) &= \str{\bs\bs}\ \mathit{esc}(u) \\
\mathit{esc}(\str{[}\,u) &= \str{\bs[}\ \mathit{esc}(u) \\
\mathit{esc}(c\,u) &= c\ \mathit{esc}(u)
  && \text{if } c \notin \{\str{\bs}, \str{[}\}
\end{aligned}
\end{equation}
Every backslash and every $\str{[}$ is escaped, so the escaped text
contains no $\Open$, and every backslash in it begins an escape, so no
$\Close$ is read from it either.  In text mode it lexes back to $u$, one
character lexeme per character of $u$, whatever follows it:
\begin{equation}\label{eq:escape-lexes}
\text{for all } u \text{ and } W,\quad
\langle \TM, i, \mathit{esc}(u)\,W \rangle \text{ reads the characters of }
u \text{ and then lexes } W \text{ as } \langle \TM, i', W \rangle
  \text{ does}
\end{equation}
where $i' = i + |\mathit{esc}(u)|$.  This is an obligation, checked by the
regressions on texts that contain $\Open$, $\Close$, and trailing
backslashes.

\subsection{The host printer}

The host printer is the inverse partner of the host parser: it prints the
statements of one bracket back to host input pieces, code with holes, where
each hole is a markup literal for the core to print:
\begin{equation}\label{eq:host-printer}
\mathcal{Q}_W(\bar\sigma) \in (\kw{pcode}(u) \mid \kw{plit}(L))^*
\end{equation}
It is given the source $W$ that the statements were parsed from, because
the statements record their comments only by span; for a span
$s = [i, j)$, $W[s]$ is the source from $i$ to $j$.  The host decides the
layout of the code, including the comments, which it should keep.  Its
requirements, (H6) and (H8), are stated with the round trip below.  The
trivial host printer copies each statement from its source, cutting the
code at its literals, $W[c_1]\,\vs \cdots W[c_m]\,\vs\ W[s]$ for a send
with comments $c_1 \cdots c_m$.

\subsection{Printing a tree}

The core prints the pieces, and the literals in their holes, which are
markup and contain brackets of their own:
\begin{equation}\label{eq:print-message}
\begin{aligned}
\mathit{show}_W(\kw{pcode}(u)\,\bar p) &= u\ \mathit{show}_W(\bar p),
  \qquad \mathit{show}_W(\epsilon) = \epsilon \\
\mathit{show}_W(\kw{plit}(L)\,\bar p) &=
  \mathit{showLit}_W(L)\ \mathit{show}_W(\bar p) \\
\mathit{showLit}_W(\kw{markup}(d, n_1 \cdots n_m, s)) &=
  \Open\ \mathit{showNode}_W(n_1) \cdots \mathit{showNode}_W(n_m)\ \Close \\
\mathit{showNode}_W(\kw{block}(u, s_u, \bar\sigma, s)) &=
  \mathit{esc}(u)\ \Open\ \mathit{show}_W(\mathcal{Q}_W(\bar\sigma))\ \Close \\
\mathit{showNode}_W(\kw{final}(u, s)) &= \mathit{esc}(u) \\
\mathit{br}_W(\sigma) &= \Open\ \mathit{show}_W(\mathcal{Q}_W(\sigma))\ \Close
\end{aligned}
\end{equation}
A message of the tree is printed in a bracket of its own.  A literal's
text is uninterpreted, so it is escaped but not trimmed or collapsed.
The children of a tree are printed in order, and the creating message of a
block before its children and its closing message after them:
\begin{equation}\label{eq:print-tree}
\begin{aligned}
\mathit{print}_W(\kw{tree}(\kappa, t_1 \cdots t_m, \bar\sigma, f)) &=
  \mathit{item}_W(t_1) \cdots \mathit{item}_W(t_m) \\
\mathit{item}_W(\kw{run}(u, s)) &= \mathit{esc}(u) \\
\mathit{item}_W(\kw{stray}(u, s)) &= \mathit{esc}(u) \\
\mathit{item}_W(\kw{note}(s)) &= \mathit{br}_W(\kw{comment}(s)) \\
\mathit{item}_W(t) &= \mathit{br}_W(\sigma)
  \quad\text{if } t \in \{\kw{mut}(\sigma), \kw{dnu}(\sigma),
    \kw{rootPop}(\sigma), \kw{refused}(\sigma)\} \\
\mathit{item}_W(\kw{sub}(\kw{tree}(\kappa, \bar t, \sigma\,\bar\sigma, f))) &=
  \mathit{br}_W(\sigma)\ \mathit{print}_W(\kw{tree}(\kappa, \bar t,
    \sigma\,\bar\sigma, f))\ \mathit{closer}_W(\bar\sigma, f) \\
\mathit{closer}_W(\bar\sigma\,\sigma', f) &= \mathit{br}_W(\sigma')
  \quad\text{if } f \in \{\kw{closed}, \kw{rejected}\} \\
\mathit{closer}_W(\epsilon, f) &= \epsilon
  \quad\text{if } f \in \{\kw{closed}, \kw{rejected}\} \\
\mathit{closer}_W(\bar\sigma, \kw{unclosed}) &= \epsilon
\end{aligned}
\end{equation}
The root is printed through its children alone: it has no creating or
closing message, and the messages that mutated it are its $\kw{mut}$
markers.  A leaf received only its creating message, so it has no closer.  A
block that was pushed and closed received at least its creating and its
closing message, and the last is the closer.  Its mutating messages are
printed from its $\kw{mut}$ markers, where they were received.  An unclosed
block has no closer; it can only be on the rightmost path of the tree, so
nothing is printed after it that it could swallow.

\paragraph{Layout.}
This printer emits no layout: it is canonical.  Whitespace rule~5 leaves the
line layout to the printer, and a pretty printer may insert white space
wherever the canonical output puts a bracket with no text before it whose
ToS is a $\kw{flow}$ block: there white space adds no child.  Layout is not
modelled further here.

\subsection{Round trip}

The round trip needs three more requirements of the host.
\begin{enumerate}
\item[(H6)] \emph{Printed statements reparse.}  If $\bar\sigma'$ is a
      nonempty sequence of statements, each a result of $\mathcal{P}$ on
      some bracket of $W$ or given by the field interface
      (Section~\ref{sec:field}), then $\mathcal{P}$ on the pieces of the
      bracket $\Open\,\mathit{show}_W(\mathcal{Q}_W(\bar\sigma'))\,\Close$,
      wherever it is placed, yields $\bar\sigma'$, up to spans.  The round
      trip needs this only for a sublist of one bracket; editing
      (Section~\ref{sec:editing}) needs the rest.
\item[(H7)] \emph{Positions are invisible.}  The host evaluator and the
      block signature are deterministic, and they relate inputs that are
      equal up to spans to outcomes that are equal up to spans.
\item[(H8)] \emph{Printing is stable.}  In the situation of (H6), if
      $W'$ contains that bracket and $\bar\sigma''$ is what $\mathcal{P}$
      yields on it, then
      $\lfloor \mathcal{Q}_{W'}(\bar\sigma'') \rfloor =
      \lfloor \mathcal{Q}_W(\bar\sigma') \rfloor$: printing what was
      printed changes nothing.
\end{enumerate}
These are what is asked of a pretty printer: that it change layout only,
keep the comments in place, and reach a fixed point after one pass.  The
trivial printer satisfies them for the statements of $W$; a statement given
by the field interface has no source to copy, and needs a printer that
prints from abstract syntax.
Write $\lfloor \cdot \rfloor$ for the canonical form: every span replaced
by $[0, 0)$, in statements, literals, and trees, and in host data as the
host defines; and every list of children $\bar t$ replaced by
$\mathit{merge}(\bar t)$, which merges adjacent runs and adjacent strays:
\begin{equation}\label{eq:merge}
\begin{aligned}
\mathit{merge}(\kw{run}(u, s)\,\kw{run}(u', s')\,\bar t) &=
  \mathit{merge}(\kw{run}(\mathit{collapse}(u\,u'), s)\,\bar t) \\
\mathit{merge}(\kw{stray}(u, s)\,\kw{stray}(u', s')\,\bar t) &=
  \mathit{merge}(\kw{stray}(u\,u', s)\,\bar t) \\
\mathit{merge}(t\,t'\,\bar t) &= t\ \mathit{merge}(t'\,\bar t)
  && (\dagger) \\
\mathit{merge}(t) &= t, \qquad \mathit{merge}(\epsilon) = \epsilon
\end{aligned}
\end{equation}
where ($\dagger$) applies unless $t$ and $t'$ are both runs or both strays.
\begin{equation}\label{eq:canonical}
\tau \approx \tau' \iff \lfloor \tau \rfloor = \lfloor \tau' \rfloor
\end{equation}
Adjacent runs arise only around a bracket that sends nothing, which leaves
no trace in the tree; the printer then prints the two texts together, and
they come back as one run.

\paragraph{Round trip (obligation).}
\begin{multline}\label{eq:round-trip}
\mathit{parse}(W) = \kw{parsed}(\bar n),\quad
\langle \chi, \kappa_{\mathrm{doc}}, \bar n \rangle \Downarrow_{\mathrm{I}}
  \kw{done}(\chi', \tau),\quad
W' = \mathit{print}_W(\tau)
\implies \\
\mathit{parse}(W') = \kw{parsed}(\bar n'),\quad
\langle \chi, \kappa_{\mathrm{doc}}, \bar n' \rangle \Downarrow_{\mathrm{I}}
  \kw{done}(\chi'', \tau'),\quad
\tau' \approx \tau,\quad
\lfloor \chi'' \rfloor = \lfloor \chi' \rfloor
\end{multline}
The reinterpretation sends the same messages, with arguments equal up to
spans, in the same order, to blocks in the same states, so by (H6) and (H7)
it has the same effects; the escaping lemma gives back the text of every
run and stray; closing a paragraph again trims nothing more, since its
runs were trimmed the first time; and closing a cell again drops nothing
more and rejects it exactly when it was rejected before.

\paragraph{Idempotence (obligation).}
Printing the reinterpreted tree reproduces the printed source exactly:
\begin{equation}\label{eq:print-idempotent}
\text{in the situation of (\ref{eq:round-trip}),}\quad
\mathit{print}_{W'}(\tau') = W'
\end{equation}
Every message in $W'$ prints again to the text it was printed from, by
(H8) and induction on its literals, and the text of every run is already
canonical.

\paragraph{Example.}
The source of the example of Section~\ref{sec:interpretation} prints as
\[
\str{[\bs{}paragraph\bs]Hi\ [\bs{}emphasis\bs]all[\bs{}endEmphasis\bs][\bs{}endParagraph\bs]}
\]
The leading space and the whitespace-only text before
$\str{[\bs{}endParagraph\bs]}$ are gone, since trimming removed them; the
space after $\str{Hi}$ is kept.  Interpreting this source gives the same
tree up to spans, and printing it again gives the same source.

%----------------------------------------------------------------------
\section{The standard library}
\label{sec:library}

The standard library is an instance of the block signature for the \NS{}
host: the block classes a document gets unless it overrides them.  This
section gives its structural core: the top-level block, sections,
paragraphs, the inline styles, lists, quotes, tables, figures, verbatim and
code displays, embedded HTML, and labels.  Citations and bibliographies,
references to labels, media, IDE links, and amplets are deferred; they need
registries and a resolution pass after interpretation.  A document that
overrides a block class extends this signature; overriding is not modelled
here.

\subsection{Values and block states}

The library inspects only a few kinds of host value: strings $\ulcorner u
\urcorner$, $\kw{nil}$, literals $\ulcorner L \urcorner$, and trees
$\tau$.  A symbol $\#x$ is a string literal: it denotes the string
$\ulcorner x \urcorner$.  Strings and trees are \emph{inline values},
written $h$: the content of a title or heading, shown as its text or as the
tree.  A message that takes inline content accepts a literal, which it
interprets with the inline root $\kappa_{\mathrm{inl}}$, or an inline value
as it is:
\begin{gather}
\label{eq:lib-inline-lit}
\rulelabel{Inline-Lit}\quad
\frac{\mathit{interpretWith}(\chi, L, \kappa_{\mathrm{inl}}) =
        \kw{done}(\chi', \tau)}
     {\langle \chi, \ulcorner L \urcorner \rangle \Downarrow_{\mathrm{H}}
        \kw{val}(\chi', \tau^{L})}
\\
\label{eq:lib-inline-raise}
\rulelabel{Inline-Raise}\quad
\frac{\mathit{interpretWith}(\chi, L, \kappa_{\mathrm{inl}}) =
        \kw{raised}(\chi', \xi)}
     {\langle \chi, \ulcorner L \urcorner \rangle \Downarrow_{\mathrm{H}}
        \kw{exn}(\chi', \xi)}
\\
\label{eq:lib-inline-value}
\rulelabel{Inline-Value}\quad
\langle \chi, h \rangle \Downarrow_{\mathrm{H}} \kw{val}(\chi, h)
\\
\label{eq:lib-inline-bad}
\rulelabel{Inline-Bad}\quad
\frac{v \text{ is neither a literal nor an inline value}}
     {\langle \chi, v \rangle \Downarrow_{\mathrm{H}} \kw{bad}(\chi)}
\end{gather}

The block states, with their content classes, are:
\begin{equation}\label{eq:lib-states}
\begin{array}{@{}ll@{}}
\kw{flow} & \kw{doc}(t, t', \bar a),\ \kw{section}(d, h, \ell),\
  \kw{list},\ \kw{numbered},\ \kw{item},\ \kw{quote}(a),\ \kw{table}(\ell), \\
  & \kw{headerRow},\ \kw{row},\ \kw{figure}(c, \ell),\
  \kw{verbatim}(u),\ \kw{codeBlock}(u, u'),\ \kw{html}(u),\
  \kw{amplet}(v) \\
\kw{para} & \kw{para} \\
\kw{cell} & \kw{cell} \\
\kw{inline} & \kw{emph},\ \kw{strong},\ \kw{code},\ \kw{link}(u),\
  \kw{inline},\ \kw{newline},\ \kw{inlineAmplet}(v)
\end{array}
\end{equation}
Here $t$, $t'$, $a$, and $c$ are an inline value or $\kw{nil}$: the title,
subtitle, attribution, and caption, each set at most once; $\bar a$ are the
authors; $\ell$ is a label, a string or $\kw{nil}$; $u'$ is the language of
a code display; $v$ is the host value of an amplet's code; and $d \ge 1$ is the
depth of a section, its level.  The top-level state is
$\kappa_{\mathrm{doc}} = \kw{doc}(\kw{nil}, \kw{nil}, \epsilon)$ and the
inline root is $\kappa_{\mathrm{inl}} = \kw{inline}$.  The verbatim, code,
HTML, amplet, and newline states are leaves: they are never pushed, so they
never receive a message.

\subsection{Content models}

The content model of a block is the set of selectors it understands, its
public protocol.  It depends only on the kind of the block, not on its
fields.  Inline content $I$, displays $D$, and flow
content $F$ are shared, as the design notes' mixins are:
\begin{equation}\label{eq:lib-models}
\begin{aligned}
I &= \{\str{emphasis}, \str{strong}, \str{code}, \str{newline},
  \str{amplet:}, \str{link:}\} \\
D &= \{\str{verbatim:}, \str{codeBlock:}, \str{codeBlock:language:},
  \str{html:}, \str{amplet:}\} \\
F &= D \cup \{\str{paragraph}, \str{list}, \str{numberedList},
  \str{quote}, \str{table}, \str{figure}\}
\end{aligned}
\end{equation}
and the content models are
\begin{equation}\label{eq:lib-model-table}
\begin{array}{@{}l@{\;}l@{}}
\kw{doc} & F \cup \{\str{section:}, \str{title:}, \str{subtitle:},
  \str{author:}\} \\
\kw{section} & F \cup \{\str{section:}, \str{label:}, \str{endSection}\} \\
\kw{para} & I \cup \{\str{endParagraph}\} \\
\kw{emph},\ \kw{strong},\ \kw{code} & I \cup \{\mathit{end}\} \\
\kw{link}(u) & (I \setminus \{\str{link:}\}) \cup \{\str{endLink}\} \\
\kw{inline} & I \\
\kw{list},\ \kw{numbered} & \{\str{item}, \mathit{end}\} \\
\kw{item} & F \cup \{\str{endItem}\} \\
\kw{quote} & F \cup \{\str{attribution:}, \str{endQuote}\} \\
\kw{table} & \{\str{headerRow}, \str{row}, \str{label:},
  \str{endTable}\} \\
\kw{headerRow},\ \kw{row} & \{\str{cell}, \str{endRow}\} \\
\kw{cell} & F \cup I \cup \{\str{endCell}\} \\
\kw{figure} & D \cup \{\str{caption:}, \str{label:}, \str{endFigure}\}
\end{array}
\end{equation}
where $\mathit{end}$ is the block's own closer, as in
$\str{endEmphasis}$ for $\kw{emph}$ and $\str{endNumberedList}$ for
$\kw{numbered}$.  Leaves understand nothing.  Each kind of block has
exactly one closer, so a closer sent to the wrong block is not understood.
A section opens only in the top-level block or a section, so its depth is
its nesting depth.

\subsection{Effects}

Every selector $x$ in a content model is one of the following.  An
\emph{opener} pushes a fresh block,
\begin{equation}\label{eq:lib-openers}
\begin{array}{@{}l@{}}
\str{paragraph} \mapsto \kw{para},\
\str{emphasis} \mapsto \kw{emph},\
\str{strong} \mapsto \kw{strong},\
\str{code} \mapsto \kw{code},\
\str{list} \mapsto \kw{list}, \\
\str{numberedList} \mapsto \kw{numbered},\
\str{item} \mapsto \kw{item},\
\str{quote} \mapsto \kw{quote}(\kw{nil}),\
\str{table} \mapsto \kw{table}(\kw{nil}), \\
\str{headerRow} \mapsto \kw{headerRow},\
\str{row} \mapsto \kw{row},\
\str{cell} \mapsto \kw{cell},\
\str{figure} \mapsto \kw{figure}(\kw{nil}, \kw{nil})
\end{array}
\end{equation}
a \emph{closer}, any selector beginning $\str{end}$, pops, and the rest
are given by $\mathit{plain}$, below, or are \emph{inline-argument}
messages, the set
\[
H = \{\str{section:}, \str{title:}, \str{subtitle:}, \str{author:},
\str{attribution:}, \str{caption:}\}.
\]  Writing $x \in \kappa$ for
$x$ in the content model of $\kappa$:
\begin{gather}
\label{eq:lib-plain}
\rulelabel{Lib-Plain}\quad
\frac{x \in \kappa \qquad x \notin H \qquad
      \mathit{plain}(\kappa, x, \bar v) = \varepsilon}
     {\langle \chi, \kappa, x, \bar v \rangle \Downarrow_{\mathrm{S}}
        \kw{eff}(\chi, \varepsilon)}
\\
\label{eq:lib-dnu}
\rulelabel{Lib-Dnu}\quad
\frac{x \notin \kappa}
     {\langle \chi, \kappa, x, \bar v \rangle \Downarrow_{\mathrm{S}}
        \kw{eff}(\chi, \kw{dnu})}
\end{gather}
\begin{equation}\label{eq:lib-plain-def}
\begin{aligned}
\mathit{plain}(\kappa, x, \epsilon) &= \kw{push}(\kappa')
  && \text{if } x \mapsto \kappa' \text{ is an opener} \\
\mathit{plain}(\kappa, x, \epsilon) &= \kw{pop}
  && \text{if } x \text{ is a closer} \\
\mathit{plain}(\kappa, \str{newline}, \epsilon) &= \kw{leaf}(\kw{newline}) \\
\mathit{plain}(\kappa, \str{verbatim:}, \ulcorner u \urcorner) &=
  \kw{leaf}(\kw{verbatim}(u)) \\
\mathit{plain}(\kappa, \str{codeBlock:}, \ulcorner u \urcorner) &=
  \kw{leaf}(\kw{codeBlock}(u, \str{newspeak})) \\
\mathit{plain}(\kappa, \str{codeBlock:language:},
  \ulcorner u \urcorner\,\ulcorner u' \urcorner) &= \kw{leaf}(\kw{codeBlock}(u, u')) \\
\mathit{plain}(\kappa, \str{html:}, \ulcorner u \urcorner) &=
  \kw{leaf}(\kw{html}(u)) \\
\mathit{plain}(\kappa, \str{label:}, \ulcorner u \urcorner) &=
  \kw{mutate}(\kappa[\ell := \ulcorner u \urcorner])
  && \text{if } \ell(\kappa) = \kw{nil} \\
\mathit{plain}(\kappa, \str{amplet:}, v) &= \mathit{amplet}(\kappa, v) \\
\mathit{plain}(\kappa, \str{link:}, \ulcorner u \urcorner) &=
  \kw{push}(\kw{link}(u))
\end{aligned}
\end{equation}
and $\mathit{plain}(\kappa, x, \bar v) = \kw{refuse}$ for every other
$x \in D \cup \{\str{label:}\}$ and $\bar v$.  So a display whose argument is not a string, and a label that is not a
string or that would relabel a block, are refused.  Since a symbol is a
string literal, $\str{label:}~\#\mathit{intro}$ and
$\str{label:}~\ulcorner\mathit{intro}\urcorner$ are the same message.

An \emph{amplet} is live content that the view computes, as in Ampleforth.
Its argument is the host value of its code: in the Newspeak host, a block
literal, such as $\str{amplet:}~[\mathit{section}\!:\ldots]$, whose value is
a closure that has not run.  Interpretation never runs it; the view does,
when it shows the amplet, and keeps the result while the amplet's code is
unchanged.  In flow content an amplet is a display; anywhere else, an inline
leaf:
\begin{equation}\label{eq:lib-amplet}
\mathit{amplet}(\kappa, v) =
\begin{cases}
\kw{leaf}(\kw{amplet}(v)) & \text{if } v \text{ is a host value and }
  \mathit{cls}(\kappa) = \kw{flow} \\
\kw{leaf}(\kw{inlineAmplet}(v)) & \text{if } v \text{ is a host value and }
  \mathit{cls}(\kappa) \ne \kw{flow} \\
\kw{refuse} & \text{otherwise}
\end{cases}
\end{equation}
A host value here is a value that is not a literal, a tree, a string, a
number, or $\kw{nil}$.

A \emph{link} is an inline block whose target $u$ is a string; its text is
inline content, but not another link, and $\str{endLink}$ closes it:
\begin{equation}\label{eq:lib-link}
\mathit{plain}(\kappa, \str{link:}, v) = \kw{refuse}
  \quad\text{if } v \text{ is not a string.}
\end{equation}  An inline-argument
message checks first that its field is free, then takes its argument as
inline content, then changes the block:
\begin{gather}
\label{eq:lib-inline}
\rulelabel{Lib-Inline}\quad
\frac{x \in \kappa \cap H \qquad \mathit{free}(\kappa, x) \qquad
      \langle \chi, v \rangle \Downarrow_{\mathrm{H}} \kw{val}(\chi', h)}
     {\langle \chi, \kappa, x, v \rangle \Downarrow_{\mathrm{S}}
        \kw{eff}(\chi', \mathit{with}(\kappa, x, h))}
\\
\label{eq:lib-inline-taken}
\rulelabel{Lib-Inline-Taken}\quad
\frac{x \in \kappa \cap H \qquad \neg\mathit{free}(\kappa, x)}
     {\langle \chi, \kappa, x, v \rangle \Downarrow_{\mathrm{S}}
        \kw{eff}(\chi, \kw{refuse})}
\\
\label{eq:lib-inline-refuse}
\rulelabel{Lib-Inline-Refuse}\quad
\frac{x \in \kappa \cap H \qquad \mathit{free}(\kappa, x) \qquad
      \langle \chi, v \rangle \Downarrow_{\mathrm{H}} \kw{bad}(\chi)}
     {\langle \chi, \kappa, x, v \rangle \Downarrow_{\mathrm{S}}
        \kw{eff}(\chi, \kw{refuse})}
\\
\label{eq:lib-inline-exn}
\rulelabel{Lib-Inline-Raise}\quad
\frac{x \in \kappa \cap H \qquad \mathit{free}(\kappa, x) \qquad
      \langle \chi, v \rangle \Downarrow_{\mathrm{H}} \kw{exn}(\chi', \xi)}
     {\langle \chi, \kappa, x, v \rangle \Downarrow_{\mathrm{S}}
        \kw{exn}(\chi', \xi)}
\end{gather}
\begin{equation}\label{eq:lib-with}
\begin{aligned}
\mathit{with}(\kappa, \str{section:}, h) &=
  \kw{push}(\kw{section}(\mathit{depth}(\kappa) + 1, h, \kw{nil})) \\
\mathit{with}(\kw{doc}(t, t', \bar a), \str{title:}, h) &=
  \kw{mutate}(\kw{doc}(h, t', \bar a)) \\
\mathit{with}(\kw{doc}(t, t', \bar a), \str{subtitle:}, h) &=
  \kw{mutate}(\kw{doc}(t, h, \bar a)) \\
\mathit{with}(\kw{doc}(t, t', \bar a), \str{author:}, h) &=
  \kw{mutate}(\kw{doc}(t, t', \bar a\,h)) \\
\mathit{with}(\kw{quote}(a), \str{attribution:}, h) &=
  \kw{mutate}(\kw{quote}(h)) \\
\mathit{with}(\kw{figure}(c, \ell), \str{caption:}, h) &=
  \kw{mutate}(\kw{figure}(h, \ell))
\end{aligned}
\end{equation}
where $\mathit{depth}$ is $0$ for $\kw{doc}$ and $d$ for
$\kw{section}(d, h, \ell)$.  A field is free when it is $\kw{nil}$:
$\mathit{free}(\kappa, x)$ holds when $x$ is $\str{title:}$,
$\str{subtitle:}$, $\str{attribution:}$, or $\str{caption:}$ and the field
it sets, $t$, $t'$, $a$, or $c$, is $\kw{nil}$, and always when $x$ is
$\str{section:}$ or $\str{author:}$.  Checking the
field before interpreting the argument means that a refused title is never
interpreted.  The library uses no $\kw{popReject}$ yet; the bibliographic
entries, whose closer checks the required fields, will.

\paragraph{What the checks cannot see.}
A check sees only the state of the block that receives the message.  Pushing
a child does not change its parent's state, so the library cannot check
the order of children, such as a title before any section, nor their
number, such as a list with no items.  Checks of that kind need either
states that count their children or a check over the finished tree.

\subsection{Errors}

A title, heading, attribution, or caption interpreted from a literal is a
tree held in a field, and its errors are the library's to report
(Section~\ref{sec:interp-results}).  It reports them with the block that
holds the field, after that block's own flag:
\begin{equation}\label{eq:lib-errors}
\begin{split}
\mathit{errors}_{\mathrm{lib}}(\kw{tree}(\kappa, \bar t, \bar\sigma, f)) ={}&
  \mathit{flag}(f)\
  \mathit{errors}_{\mathrm{lib}}(\tau_1) \cdots
  \mathit{errors}_{\mathrm{lib}}(\tau_k) \\
  & \mathit{errors}_{\mathrm{lib}}(t_1) \cdots \mathit{errors}_{\mathrm{lib}}(t_m)
\end{split}
\end{equation}
where $\tau_1 \cdots \tau_k$ are the trees among the fields of $\kappa$, in
order, and $\mathit{errors}_{\mathrm{lib}}$ is $\mathit{errors}$ on every
child that is not a block.

\paragraph{Example.}
The table of the design notes,
\[
\begin{array}{@{}l@{}}
\str{[\bs{}table.\ headerRow.\ cell\bs]Act[\bs{}endCell.\ cell\bs]Scenes[\bs{}endCell.\ endRow\bs]} \\
\str{[\bs{}row.\ cell\bs]I[\bs{}endCell.\ cell\bs]5[\bs{}endCell.\ endRow\bs]} \\
\str{[\bs{}endTable\bs]}
\end{array}
\]
interprets to a $\kw{table}(\kw{nil})$ holding a $\kw{headerRow}$ and a
$\kw{row}$, each holding two cells with one run each: the line breaks
between the brackets fall in the table, a $\kw{flow}$ block, and add
nothing.  A cell holding $\str{x[\bs{}paragraph\bs]y[\bs{}endParagraph.\
endCell\bs]}$ mixes text with a paragraph and is rejected when it closes;
$\str{[\bs{}cell.\ paragraph\bs]y[\bs{}endParagraph.\ endCell\bs]}$ is a
cell holding flow content, and $\str{[\bs{}cell\bs]\ y\
[\bs{}endCell\bs]}$ is the shorthand, a cell holding the run $\str{y}$.

%----------------------------------------------------------------------
\section{Editing}
\label{sec:editing}

A live view shows the tree of a document and lets its reader edit it.  In a
mark-edit language every such edit is carried back to source: the design
notes require that a live-view edit change source, never block instances.
This section gives the edit layer.  An edit never changes a tree.  It
rewrites a source, and the new tree is the interpretation of the new
source; what the edit layer promises is that this tree is the one the
reader asked for.

The layer has four steps.  A \emph{point} is a position in a finished tree,
which is how the view reports where an edit happens.  Its \emph{anchor} is a
position among the items of the source (Section~\ref{sec:interpretation}),
the place where interpretation would add whatever is put there to the right
block.  A \emph{splice} replaces the items between two anchors with new
items and prints them into the source.  Each edit of the view is a splice,
chosen so that the reinterpreted source shows the edit, and tidied so that
it leaves no empty markup behind.

\paragraph{The running example.}  Most examples below edit
\[
W_1 = \str{[\bs{}paragraph\bs]Hi\ [\bs{}emphasis\bs]all[\bs{}endEmphasis.\
  endParagraph\bs]}
\]
Its nodes are $n_0$, with empty contents and the statement
$\str{paragraph}$; $n_1$, with the contents $\str{Hi}\vs$ and the statement
$\str{emphasis}$; and $n_2$, with the contents $\str{all}$ and the
statements $\str{endEmphasis}$ and $\str{endParagraph}$.  The root of its
tree has one child, the paragraph $P$, whose children are
$\kw{run}(\str{Hi}\vs, s_1)$ and the emphasis $E$; the one child of $E$ is
$\kw{run}(\str{all}, s_2)$, for $s_1, s_2$ the spans of the contents of
$n_1$ and $n_2$.  Every example in this section is a regression of the
library suite, named in brackets, so what it claims is checked against the
SpecTec transcription.

\subsection{The image and provenance}
\label{sec:image}

The design notes give text in the view one of three provenances: a literal
in this document; a literal elsewhere in the image, such as in a method
that serves as a template, whose edit goes to that method's source; or a
computation, which is read-only.

An \emph{image} $I$ is a finite map from source names $d$ to sources.  It
holds the document's source, $I(d_0)$, and the sources of the host's
methods.  A literal records its source, $\kw{markup}(d, \bar n, s)$, and a
tree held in a value records its literal, $\tau^{L}$
(Section~\ref{sec:interpretation}).  A \emph{markup} $m$ is either the
document, $\kw{doc}(d)$, or a literal; $\mathit{src}(m)$ is its source name
and $\mathit{nodes}(m)$ its DocumentBlockASTs.  The tree of the document is
the interpretation of $\kw{doc}(d_0)$, and the tree $\tau^L$ held in a
field is that of $L$.  The provenance of a markup is
\begin{equation}\label{eq:provenance}
\mathit{prov}(m) =
\begin{cases}
\kw{local} & \text{if } \mathit{src}(m) = d_0 \\
\kw{template} & \text{if } \mathit{src}(m) \in \mathrm{dom}(I)
  \setminus \{d_0\} \\
\kw{computed} & \text{otherwise}
\end{cases}
\end{equation}
A literal parsed at run time from a string has a source outside the image,
so its tree is computed; so is a value that is no tree at all, such as a
string that code computed and a block shows as its title.

Two requirements of the host make provenance trustworthy:
\begin{enumerate}
\item[(H9)] \emph{Literals are current.}  If $\mathit{src}(L) = d \in
      \mathrm{dom}(I)$, the text of $I(d)$ at the span of $L$ is the text
      that $L$ was parsed from.
\item[(H10)] \emph{A literal can be spliced.}  If a source of the host
      holds a literal, replacing the text inside the literal's brackets by
      text that parses as markup gives a source that the host parses as
      before except for that literal, and evaluation uses the new literal
      from then on.
\end{enumerate}

\paragraph{Example.}  Each of these documents has a title that is a tree,
of each provenance in turn:
\[
\begin{array}{@{}ll@{}}
W_{\mathrm{l}} = \str{[\bs{}title:\ [\bs{}Hamlet\bs]\bs][\bs{}paragraph\bs]x[\bs{}endParagraph\bs]}
  & \kw{local} \\
W_{\mathrm{t}} = \str{[\bs{}title:\ standardTitle\bs]} & \kw{template} \\
W_{\mathrm{c}} = \str{[\bs{}title:\ scratchTitle\bs]} & \kw{computed}
\end{array}
\]
In $W_{\mathrm{l}}$ the literal is the document's own.  In
$W_{\mathrm{t}}$ the host method $\str{standardTitle}$ answers the literal
$\str{\^{}[\bs{}The\ Play\bs]}$ of the source $\str{Templates}$, which the
image holds, and typing $\str{Great}\vs$ before $\str{Play}$ rewrites that
source, not the document [edit-template].  In $W_{\mathrm{c}}$ the method
$\str{scratchTitle}$ answers a literal parsed at run time, whose source is
not in the image, and typing in its title is refused [edit-computed-tree].

\subsection{The field interface}
\label{sec:field}

Besides the tree, the view needs to know which fields a block displays, and
where each came from.  This is the field interface.  For a block state
$\kappa$ and the messages $\bar\sigma$ the block received, it gives the
block's \emph{fields}
\begin{equation}\label{eq:fields}
\mathit{fields}(\kappa, \bar\sigma) = \varphi_0 \cdots \varphi_{r-1}
\qquad
\varphi ::= \langle v, \sigma, k \rangle
\end{equation}
each a value $v$ that the block displays, such as its title, together with
the message $\sigma \in \bar\sigma$ and the index $k$ of the argument it
came from.  It also says which blocks an edit may merge,
$\mathit{merge}(\kappa)$, and gives, for each style $x$, the statements
that open and close it, $\mathit{open}(x)$ and $\mathit{close}(x)$.  The
host gives, for an argument expression $\eta$: its span
$\mathit{span}(\eta)$; $\mathit{str}(\eta) = s$ when $\eta$ is a string
literal standing alone, at span $s$; $\mathit{enc}(u)$, the text of a
string literal whose value is the string $u$; and, where it can,
$\mathcal{Q}_{\mathrm{v}}(v)$, the text of an expression whose value is $v$.
\begin{enumerate}
\item[(H11)] \emph{Strings can be edited.}  Replacing the text at the span
      of a string literal $\eta$ by $\mathit{enc}(u)$ gives a source that
      parses as before, except that $\eta$ now evaluates, without effect,
      to the string $\ulcorner u \urcorner$.
\item[(H12)] \emph{Detaching keeps the value.}  Replacing the text at
      $\mathit{span}(\eta)$ by $\mathcal{Q}_{\mathrm{v}}(v)$ gives a source
      that parses as before, except that $\eta$ now evaluates, without
      effect, to $v$ up to $\approx$; a tree $\tau^{L}$ is printed as a copy
      of $L$, which the method receiving it interprets to $\tau$ again.
\end{enumerate}
The standard library's field interface is given in
Section~\ref{sec:lib-field}.

\paragraph{Example.}  Let $L = \str{[\bs{}Hamlet\bs]}$ be the literal of
$W_{\mathrm{l}}$, with the tree $\tau_{\mathrm{H}}$.  The root of
$W_{\mathrm{l}}$ has the state $\kw{doc}(\tau_{\mathrm{H}}^{L}, \kw{nil},
\epsilon)$ and one field, $\langle \tau_{\mathrm{H}}^{L}, \sigma, 0
\rangle$ for $\sigma$ the $\str{title:}$ send; a subtitle of $\kw{nil}$ is
not shown.  The one field of the code block of
$\str{[\bs{}codeBlock:\ 'x\ :=\ 1'\bs]}$ is a string, from an argument
that is a string literal standing alone, so it can be set in place
[edit-string].

\subsection{Points}
\label{sec:points}

\begin{equation}\label{eq:points}
\pi ::= \epsilon \mid \pi.i \mid \pi.\varphi_j
\qquad
p ::= \kw{gap}(\pi, i) \mid \kw{in}(\pi, i, k)
\end{equation}
A path leads from the root of a tree to a block: $\pi.i$ to the $i$th child
of the block at $\pi$, and $\pi.\varphi_j$ to the root of the tree held in
its $j$th field, counting from $0$.  For the tree $\tau$ of a markup $m$,
the block $\tau@\pi$ at a path and the markup $\mathit{mk}(\pi)$ of the
tree the path enters last are
\begin{equation}\label{eq:path}
\begin{aligned}
\tau@\epsilon &= \tau & \mathit{mk}(\epsilon) &= m \\
\tau@\pi.i &= \tau' & \mathit{mk}(\pi.i) &= \mathit{mk}(\pi)
  &&\text{if } \tau@\pi = \kw{tree}(\kappa, \bar g, \bar\sigma, f),\
    g_i = \kw{sub}(\tau') \\
\tau@\pi.\varphi_j &= \tau' & \mathit{mk}(\pi.\varphi_j) &= L
  &&\text{if } \tau@\pi = \kw{tree}(\kappa, \bar g, \bar\sigma, f),\
    \mathit{fields}(\kappa, \bar\sigma)_j = \langle \tau'^{L}, \sigma, k
    \rangle
\end{aligned}
\end{equation}
A point $\kw{gap}(\pi, i)$ lies before the $i$th child of the block at
$\pi$, or at its end, and $\kw{in}(\pi, i, k)$ inside the text of the $i$th
child, before its $k$th character.  For $\tau@\pi = \kw{tree}(\kappa, \bar
g, \bar\sigma, f)$ the points are
\begin{equation}\label{eq:point-ok}
\kw{gap}(\pi, i) \ \text{ for } 0 \le i \le |\bar g|
\qquad
\kw{in}(\pi, i, k) \ \text{ for } g_i \in \{\kw{run}(u, s),
  \kw{stray}(u, s)\} \text{ and } 0 < k < |u|
\end{equation}
so the ends of a text are gaps, and no point lies inside an error leaf
other than stray text: a message that was not understood, refused, or that
would pop the root is shown as its message, and is edited through the
source.  The tree of a point is the one its path enters last,
$\mathit{mk}(p)$ is the markup of its path, and points are ordered in
document order.

For a block $B = \kw{tree}(\kappa, \bar g, \sigma_0 \cdots \sigma_r, f)$
other than a root, its \emph{creating} and \emph{closing} messages are
\begin{equation}\label{eq:op-cl}
\mathit{op}(B) = \sigma_0
\qquad
\mathit{cl}(B) = \sigma_r \ \text{ if } r \ge 1 \text{ and } f \neq
  \kw{unclosed}
\end{equation}
and $\mathit{cls}(B)$ is the class of $\kappa$.  An unclosed block has no
closing message.

\paragraph{Example.}  In $W_1$, $\kw{gap}(\epsilon, 1)$ is the end of the
document, $\kw{gap}(0, 0)$ the start of the paragraph, $\kw{gap}(0.1, 1)$
the end of the emphasis, and $\kw{in}(0, 0, 2)$ lies between $\str{Hi}$ and
the space [ex-point-in].  The end of the paragraph is $\kw{gap}(0, 2)$
[ex-point-gap-last]; $\kw{gap}(0, 3)$ is no point [ex-point-gap-past], and
nor is $\kw{in}(0, 0, 3)$, since the end of $\str{Hi}\vs$ is the gap
$\kw{gap}(0, 1)$ [ex-point-in-end].  In $W_{\mathrm{l}}$,
$\kw{in}(\varphi_0, 0, 3)$ lies inside the title, between $\str{Ham}$ and
$\str{let}$, and its markup is the literal [edit-title].

\subsection{Located text}

A white-space run is one character of the view but may be many of the
source, and a character of the view may be an escape.  For contents $u$,
$\mathit{col}(u) = \mathit{col}(u, 0)$ is $\mathit{collapse}(u)$ with each
character paired with the range of $u$ it came from:
\begin{equation}\label{eq:col}
\begin{aligned}
\mathit{col}(\epsilon, a) &= \epsilon \\
\mathit{col}(c\,u, a) &= \langle c, [a, a{+}1) \rangle\ \mathit{col}(u, a{+}1)
  && \text{if } c \notin \mathit{WS} \\
\mathit{col}(c\,w\,u, a) &= \langle \vs, [a, b) \rangle\ \mathit{col}(u, b)
  && \text{if } c \in \mathit{WS},\ w \in \mathit{WS}^*,\
     u \notin \mathit{WS}\,\Sigma^*,\ b = a + 1 + |w|
\end{aligned}
\end{equation}
A child $\kw{run}(u', s)$ of a tree of markup $m$ comes from the node of
$m$ whose contents $u$ have the span $s$, and $u'$ is
$\mathit{collapse}(u)$ less any space that trimming removed at either end
(\ref{eq:trim}).  The located text of a child drops the same characters:
\begin{equation}\label{eq:loc}
\begin{gathered}
\mathit{loc}(\kw{run}(u', s)) =
\begin{cases}
\mathit{col}(u) & \text{if } \mathit{collapse}(u) = u' \\
\bar r & \text{if } \mathit{collapse}(u) = \vs\,u',\ \mathit{col}(u) =
  r\,\bar r \\
\bar r & \text{if } \mathit{collapse}(u) = u'\,\vs,\ \mathit{col}(u) =
  \bar r\,r \\
\bar r & \text{if } \mathit{collapse}(u) = \vs\,u'\,\vs,\
  \mathit{col}(u) = r\,\bar r\,r'
\end{cases}
\\
\mathit{loc}(\kw{stray}(u, s)) = \langle u_0, [0, 1) \rangle \cdots
  \langle u_{|u|-1}, [|u|{-}1, |u|) \rangle
\end{gathered}
\end{equation}
So the $k$th character of the text of a child has a range $[a_k, b_k)$ in
the contents of its node, and anchors are offsets into contents with
escapes resolved: none falls inside an escape or inside a run of white
space.

\paragraph{Example.}  The source
\[
W_{\mathrm{w}} = \str{[\bs{}paragraph\bs]\ \ a\ \bs[\ \ b\ [\bs{}endParagraph\bs]}
\]
has the contents $u = \vs\vs\str{a}\vs\str{[}\vs\vs\str{b}\vs$ in its second node,
with the escape resolved, and its paragraph has the one child
$\kw{run}(\str{a}\vs\str{[}\vs\str{b}, s)$, trimmed at both ends.  The
ranges of its five characters are
\[
[2, 3)\quad [3, 4)\quad [4, 5)\quad [5, 7)\quad [7, 8)
\]
[ex-loc]: the escape $\str{\bs[}$ is one character of $u$, and the two
spaces before $\str{b}$ one character of the view.  So $\kw{in}(0, 0, 3)$,
before that space, is at offset $5$ of $u$, and $\kw{in}(0, 0, 4)$, before
$\str{b}$, at offset $7$ [ex-loc-anchor, ex-loc-anchor-b].

\subsection{Anchors}

An anchor is a position in the items of a markup, $\mathit{items}(m) =
\mathit{items}(\mathit{nodes}(m))$:
\begin{equation}\label{eq:anchors}
a ::= \kw{txt}(n, o) \mid \kw{cut}(n, k) \mid \kw{end}(m)
\end{equation}
$\kw{txt}(n, o)$ is the offset $o$ into the contents $u$ of the node $n$,
$0 \le o \le |u|$; $\kw{cut}(n, k)$ lies between the statements $k-1$ and
$k$ of the bracket of $n$, $0 < k < |\bar\sigma|$; and $\kw{end}(m)$ lies
after the last node of $m$, and is used only when that node has a bracket.
For $\mathit{nodes}(m) = n_0 \cdots n_{l}$, where $n_i$ has contents $u_i$
and statements $\bar\sigma_i$, the items before an anchor are
\begin{equation}\label{eq:anchor-items}
\begin{aligned}
\mathit{before}(\kw{txt}(n_i, o)) &= \mathit{items}(n_0 \cdots n_{i-1})\
  \kw{text}(u_i[0, o)) \\
\mathit{before}(\kw{cut}(n_i, k)) &= \mathit{items}(n_0 \cdots n_{i-1})\
  \kw{text}(u_i)\ \bar\sigma_i[0, k) \\
\mathit{before}(\kw{end}(m)) &= \mathit{items}(m)
\end{aligned}
\end{equation}
and $\mathit{after}(a)$ is the rest, so that $\mathit{items}(m) =
\mathit{before}(a) \oplus \mathit{after}(a)$, where $\oplus$ concatenates,
joining a text at the end of its left operand to a text at the start of its
right one.  For $a \le a'$, $\mathit{items}(a, a')$ is the $\bar\omega$
with $\mathit{before}(a') = \mathit{before}(a) \oplus \bar\omega$.

A point's anchor is found from the messages around it.  For $\sigma$ the
$k$th statement of $n_i$, the anchors just after and just before it are
\begin{equation}\label{eq:anchor-after}
\mathit{after}(\sigma) =
\begin{cases}
\kw{cut}(n_i, k{+}1) & \text{if } k + 1 < |\bar\sigma_i| \\
\kw{txt}(n_{i+1}, 0) & \text{if } k + 1 = |\bar\sigma_i| \text{ and } i < l \\
\kw{end}(m) & \text{if } k + 1 = |\bar\sigma_i| \text{ and } i = l
\end{cases}
\end{equation}
\begin{equation}\label{eq:anchor-before}
\mathit{before}(\sigma) =
\begin{cases}
\kw{cut}(n_i, k) & \text{if } k > 0 \\
\kw{txt}(n_i, |u_i|) & \text{if } k = 0
\end{cases}
\end{equation}
and $\mathit{start}(m)$ is $\kw{txt}(n_0, 0)$, or $\kw{end}(m)$ when $m$ has
no nodes.  The anchor after a child $g$ of a tree is
\begin{equation}\label{eq:anchor-end}
\mathit{end}(g) =
\begin{cases}
\kw{txt}(n, b) & \text{if } g \text{ is a run or stray text from } n
  \text{ whose last range is } [a, b) \\
\mathit{after}(\mathit{cl}(B)) & \text{if } g = \kw{sub}(B) \\
\mathit{after}(\sigma) & \text{if } g \in \{\kw{mut}(\sigma),
  \kw{dnu}(\sigma), \kw{refused}(\sigma), \kw{rootPop}(\sigma)\} \\
\mathit{after}(\kw{comment}(s)) & \text{if } g = \kw{note}(s)
\end{cases}
\end{equation}
and, for $\tau@\pi = \kw{tree}(\kappa, g_0 \cdots g_{r-1}, \bar\sigma, f)$,
\begin{equation}\label{eq:anchor-point}
\begin{aligned}
\mathit{anchor}(\kw{in}(\pi, i, k)) &= \kw{txt}(n, a_k)
  && \text{for } n \text{ the node of } g_i \text{ and } [a_k, b_k)
     \text{ the } k\text{th range of } \mathit{loc}(g_i) \\
\mathit{anchor}(\kw{gap}(\pi, 0)) &= \mathit{after}(\mathit{op}(\tau@\pi))
  && \text{if } \tau@\pi \text{ is not a root} \\
\mathit{anchor}(\kw{gap}(\pi, 0)) &= \mathit{start}(\mathit{mk}(\pi))
  && \text{if } \tau@\pi \text{ is a root} \\
\mathit{anchor}(\kw{gap}(\pi, i)) &= \mathit{end}(g_{i-1})
  && \text{if } i > 0
\end{aligned}
\end{equation}
An unclosed block has no closing message, so there is no point after it;
it is on the rightmost path of the tree, and the points inside it are the
ones that matter.  A gap after a run lies before the space that trimming
removed from it, if any; a caret (Section~\ref{sec:carets}) can lie after
that space.

\paragraph{Example.}  The anchors of points of $W_1$:
\[
\begin{array}{@{}lll@{}}
\kw{gap}(0, 0) & \mathit{after}(\str{paragraph}) = \kw{txt}(n_1, 0)
  & \text{[ex-anchor-first-gap]} \\
\kw{in}(0, 0, 2) & \kw{txt}(n_1, 2) & \text{[ex-anchor-in]} \\
\kw{gap}(0, 1) & \mathit{end}(\kw{run}(\str{Hi}\vs, s_1)) =
  \kw{txt}(n_1, 3) & \text{[ex-anchor-before-emph]} \\
\kw{gap}(0.1, 0) & \mathit{after}(\str{emphasis}) = \kw{txt}(n_2, 0)
  & \text{[ex-anchor-emph-start]} \\
\kw{gap}(0.1, 1) & \mathit{end}(\kw{run}(\str{all}, s_2)) =
  \kw{txt}(n_2, 3) & \text{[ex-anchor-emph-end]} \\
\kw{gap}(0, 2) & \mathit{end}(\kw{sub}(E)) =
  \mathit{after}(\str{endEmphasis}) = \kw{cut}(n_2, 1)
  & \text{[ex-anchor-after-emph]} \\
\kw{gap}(\epsilon, 1) & \mathit{end}(\kw{sub}(P)) =
  \mathit{after}(\str{endParagraph}) = \kw{end}(\kw{doc}(d_0))
  & \text{[ex-anchor-end]}
\end{array}
\]
The gaps just before the emphasis and just inside it, $\kw{gap}(0, 1)$ and
$\kw{gap}(0.1, 0)$, show the same place in the view but have different
anchors: $\mathit{items}(\kw{txt}(n_1, 3), \kw{txt}(n_2, 0))$ is the one
statement $\str{emphasis}$, so text typed at the first goes to the
paragraph and text typed at the second to the emphasis.

\subsection{Carets}
\label{sec:carets}

The view shows a point, but a point cannot always say where the next
character typed goes.  In a closed paragraph trimming removes the space at
the end of $\str{Hello}\vs$, so the last gap of the run anchors before that
space, and a character typed at it would go before the space too.  A
\emph{caret} is a point with an anchor:
\begin{equation}\label{eq:carets}
c ::= \kw{exact}(p) \mid \kw{at}(p, a)
\qquad
\mathit{pt}(c) = p
\qquad
\mathit{anchor}(\kw{exact}(p)) = \mathit{anchor}(p)
\qquad
\mathit{anchor}(\kw{at}(p, a)) = a
\end{equation}
The point is where the view draws the caret and which blocks it is in; the
anchor is where text goes.  Where an edit or a function expects a caret, a
point $p$ stands for $\kw{exact}(p)$.

An offset $k$ of the source $W$ of a markup $m$ with the nodes $n_0 \cdots
n_l$ is the anchor
\begin{equation}\label{eq:caret-anchor}
\mathit{at}_m(k) =
\begin{cases}
\kw{txt}(n_i, |W[i_u, k)|_{\mathrm{lex}}) & \text{for the least } i
  \text{ with } i_u \le k \le j_u \\
\kw{end}(m) & \text{if there is no such } i
\end{cases}
\end{equation}
where $[i_u, j_u)$ is the span of the contents of $n_i$, and
$|w|_{\mathrm{lex}}$ counts the characters $w$ lexes to in text mode, each
of the escapes $\str{\bs\bs}$ and $\str{\bs[}$ being one.

An anchor of $m$ is placed at a point of the tree $\tau$ of $m$.  If the
contents of $n_i$ gave the child $g_j$ of $\tau@\pi$, a run or stray text
whose $\mathit{loc}$ ranges are $[a_0, b_0) \cdots [a_{r-1}, b_{r-1})$,
then
\begin{equation}\label{eq:caret-resolve}
\mathit{place}(\kw{txt}(n_i, o)) =
\begin{cases}
\kw{gap}(\pi, j) & \text{if } v = 0 \\
\kw{in}(\pi, j, v) & \text{if } 0 < v < r \\
\kw{gap}(\pi, j{+}1) & \text{if } v = r
\end{cases}
\qquad
v = |\{\, t \mid b_t \le o \,\}|
\end{equation}
which is before the first character whose source ends after $o$.  If the
contents of $n_i$ gave no child, being empty or white space that was
dropped or trimmed away, then $\mathit{place}(\kw{txt}(n_i, o)) =
\mathit{after}_\tau(i)$, and $\mathit{place}(\kw{end}(m)) =
\mathit{after}_\tau(l{+}1)$, where $\mathit{after}_\tau(i)$ is the gap just
after the statements before $n_i$: $\kw{gap}(\epsilon, 0)$ for $i = 0$;
$\mathit{place}(\kw{txt}(n_{i-1}, |u_{i-1}|))$ if the bracket of $n_{i-1}$
is empty; and otherwise, for $\sigma$ its last statement and $g_j$ the
child of $\tau@\pi$ where $\sigma$ left its mark,
\[
\mathit{after}_\tau(i) =
\begin{cases}
\kw{gap}(\pi.j, 0) & \text{if } g_j = \kw{sub}(B),\ \sigma = \mathit{op}(B),
  \text{ and } B \text{ was pushed} \\
\kw{gap}(\pi, j{+}1) & \text{if } g_j = \kw{sub}(B) \text{ and } \sigma =
  \mathit{cl}(B) \\
\kw{gap}(\pi, j{+}1) & \text{if } g_j = \kw{sub}(B),\ \sigma =
  \mathit{op}(B), \text{ and } B \text{ is a leaf} \\
\kw{gap}(\pi, j{+}1) & \text{if } g_j \in \{\kw{mut}(\sigma),
  \kw{dnu}(\sigma), \kw{refused}(\sigma), \kw{rootPop}(\sigma)\} \\
\kw{gap}(\pi, j{+}1) & \text{if } \sigma = \kw{comment}(s) \text{ and } g_j =
  \kw{note}(s)
\end{cases}
\]
A block was pushed if it is unclosed or received at least two messages;
otherwise it is a leaf.  For $\pi_m$ the path to the field whose literal is
$m$, or $\epsilon$ for the document, and $\pi_m.p$ the point $p$ of $\tau$
as a point of the document, a caret is well formed when
\begin{equation}\label{eq:caret-ok}
\mathit{ok}(\kw{exact}(p))
\qquad
\mathit{ok}(\kw{at}(p, a)) \iff a \neq \kw{cut}(n, k) \wedge
  \pi_m.\mathit{place}(a) = p
\end{equation}
So the anchor of a caret that is not exact lies in white space that the
view does not show, after its point's anchor and before the next character
the view shows.

\paragraph{Examples.}  In
\[
W_{\mathrm{h}} = \str{[\bs{}paragraph\bs]Hello}\vs\str{[\bs{}endParagraph\bs]}
\]
the run is $\str{Hello}$, and $\kw{at}(\kw{gap}(0, 1), \kw{txt}(n_1, 6))$ is a caret
after the trimmed space [caret-ok-slack].  The point $\kw{in}(0, 0, 1)$
with the anchor $\kw{txt}(n_1, 2)$ is not a caret, since the $\str{e}$
between them is shown [caret-ok-shown], and neither is a caret at a cut
[caret-ok-cut].  The end of
\[
\str{[\bs{}paragraph\bs]ab[\bs{}endParagraph\bs]}
\]
is placed at the gap after the paragraph, $\kw{gap}(\epsilon, 1)$ [caret-place-end], and so is
an offset inside the two spaces dropped between
$\str{[\bs{}paragraph\bs]a[\bs{}endParagraph\bs]}$ and a second paragraph
[caret-place-dropped].  In
$\str{[\bs{}paragraph\bs]}\vs\str{there[\bs{}endParagraph\bs]}$, the
offset $1$, after the trimmed space, is placed at $\kw{gap}(0, 0)$, and the offset
$2$, after the $\str{t}$, at $\kw{in}(0, 0, 1)$ [caret-place-lead,
caret-place-in].

\subsection{Splices}

A splice $\mathit{splice}_m(a, a', \bar\omega, \bar\omega')$, for anchors
$a \le a'$ of $m$, rewrites $W = I(\mathit{src}(m))$, writing the items
$\bar\omega$ before the caret and $\bar\omega'$ after it.  For $n = \kw{block}(u, s_u,
\bar\sigma, s)$ or $\kw{final}(u, s_u)$, with $s_u = [i_u, j_u)$ and $s =
[i, j)$, the source replaced runs from $\mathit{L}(a)$ to $\mathit{R}(a')$:
\begin{equation}\label{eq:splice-bounds}
\begin{aligned}
\mathit{L}(\kw{txt}(n, o)) &= i_u &
\mathit{R}(\kw{txt}(n, o)) &= j_u \\
\mathit{L}(\kw{cut}(n, k)) &= j_u &
\mathit{R}(\kw{cut}(n, k)) &= j \\
\mathit{L}(\kw{end}(m)) &= e_m &
\mathit{R}(\kw{end}(m)) &= e_m
\end{aligned}
\end{equation}
where $e_m$, the end of the contents of $m$, is $|W|$ for a document and
the start of its closing $\Close$ for a literal; the end of the contents
of a node with a bracket is where its $\Open$ starts.  The replacement is
\begin{equation}\label{eq:splice}
\begin{aligned}
\mathit{left}(\kw{txt}(n, o)) &= \mathit{esc}(u[0, o)) &
\mathit{right}(\kw{txt}(n, o)) &= \mathit{esc}(u[o, |u|)) \\
\mathit{left}(\kw{cut}(n, k)) &=
  \Open\,\mathit{show}_W(\mathcal{Q}_W(\bar\sigma[0, k)))\,\Close &
\mathit{right}(\kw{cut}(n, k)) &=
  \Open\,\mathit{show}_W(\mathcal{Q}_W(\bar\sigma[k, |\bar\sigma|)))\,\Close \\
\mathit{left}(\kw{end}(m)) &= \epsilon &
\mathit{right}(\kw{end}(m)) &= \epsilon
\end{aligned}
\end{equation}
\begin{equation}\label{eq:out}
\mathit{out}(\epsilon) = \epsilon \qquad
\mathit{out}(\kw{text}(u)\,\bar\omega) = \mathit{esc}(u)\ \mathit{out}(\bar\omega)
\qquad
\mathit{out}(\bar\sigma\,\bar\omega) =
  \Open\,\mathit{show}_W(\mathcal{Q}_W(\bar\sigma))\,\Close\ \mathit{out}(\bar\omega)
\end{equation}
with $\bar\sigma$ a maximal nonempty sequence of statements, and
\begin{equation}\label{eq:splice-def}
\mathit{splice}_m(a, a', \bar\omega, \bar\omega') =
\begin{cases}
\kw{unchanged} & \text{if } a = a' \text{ and } \bar\omega\,\bar\omega' =
  \epsilon \\
\langle W', k \rangle & \text{otherwise}
\end{cases}
\end{equation}
\[
W' = W[0, \mathit{L}(a))\ \mathit{left}(a)\ \mathit{out}(\bar\omega)\
  \mathit{out}(\bar\omega')\ \mathit{right}(a')\ W[\mathit{R}(a'), |W|)
\qquad
k = \mathit{L}(a) + |\mathit{left}(a)| + |\mathit{out}(\bar\omega)|
\]
where $k$ is the offset of the caret in the new source.  The two sequences
are printed apart, so where the caret falls between two statements their
bracket is split, and the caret is always at a text position.  The text at
either end of a splice is printed again, whole, from the resolved
contents.  A source may hold a lone backslash or bracket that
lexes as a character only because of what follows it; $\mathit{esc}$
writes each as an escape, so nothing written next to it can change how it
lexes.  A bracket cut by an anchor is printed again as two.

\paragraph{Splicing (obligation).}
Write $\bar\omega \doteq \bar\omega'$ when the two are equal up to spans
once empty texts are removed; an empty text adds nothing
(Step-Text-None).  If $\langle W', k \rangle$ is the result of
$\mathit{splice}_m(a, a', \bar\omega, \bar\omega')$, then $W'$ parses, the
markup $m'$ in its place has
\begin{equation}\label{eq:splice-items}
\mathit{items}(m') \doteq
  \mathit{before}(a) \oplus \bar\omega\,\bar\omega' \oplus \mathit{after}(a')
\end{equation}
and every other part of the parse is as before; for a literal, by (H10).
This follows from the escaping lemma (\ref{eq:escape-lexes}) and (H6).
Since interpretation is deterministic, the interpretation of $W'$ also
reaches the same configuration as that of $W$, up to spans, just before the
item that holds $a$: every block closed before the edit is unchanged.

\paragraph{Examples.}  In
\[
W_{\mathrm{p}} = \str{[\bs{}paragraph\bs]ab[\bs{}endParagraph.\
paragraph\bs]cd[\bs{}endParagraph\bs]}
\]
the gap between the paragraphs,
$\kw{gap}(\epsilon, 1)$, has the anchor $\mathit{after}(\str{endParagraph})
= \kw{cut}(n_1, 1)$, inside the bracket of $n_1$ [ex-cut-anchor].
Splicing the statements $\str{paragraph}\ \str{endParagraph}$ there prints
that bracket again as two, around the new one [ex-cut]:
\[
\begin{array}{@{}l@{}}
\str{[\bs{}paragraph\bs]ab[\bs{}endParagraph\bs]} \\
\str{[\bs{}paragraph.\ endParagraph\bs][\bs{}paragraph\bs]cd[\bs{}endParagraph\bs]}
\end{array}
\]
At a $\kw{txt}$ anchor the contents are printed again with every escape.
Splicing $\kw{text}(\str{x})$ at offset $5$ of the contents $u$ of the
source $W_{\mathrm{w}}$ replaces their text,
$\vs\vs\str{a}\vs\str{\bs[}\vs\vs\str{b}\vs$, by
$\mathit{esc}(u[0, 5))\,\str{x}\,\mathit{esc}(u[5, 9)) =
\vs\vs\str{a}\vs\str{\bs[x}\vs\vs\str{b}\vs$, and the paragraph shows
$\str{a}\vs\str{[x}\vs\str{b}$ [ex-loc-type, edit-escape,
edit-escape-left].

\subsection{Tidying}

Closing a block where it was just opened would leave an empty inline
block, which renders nothing, so the items of an edit are tidied before
they are written.  The items are \emph{marked}: a marked item
$\sigma^{\uparrow B}_{o}$ opens the block $B$ and $\sigma^{\downarrow
B}_{o}$ closes it, where the origin $o$ is $\mathrm{l}$, $\mathrm{s}$, or
$\mathrm{r}$ as the item lies before the splice, in it, or after it; a new
style $x$ counts as a block identified by $\mathit{open}(x)$.  Each
$\mathit{op}(B)$ and $\mathit{open}(x)$ that an edit writes is marked
$\cdot^{\uparrow}_{\mathrm{s}}$, each $\mathit{cl}(B)$ and
$\mathit{close}(x)$ is marked $\cdot^{\downarrow}_{\mathrm{s}}$, and so, in
the items between the anchors that styling writes again, are the closing
messages of the $L_i$ and the creating messages of the $R_i$
(Section~\ref{sec:edits}).  For points $p \le q$, with blocks $B_0 \cdots
B_h$ and $B'_0 \cdots B'_{h'}$ on their paths, the marked messages just
outside the splice are
\begin{equation}\label{eq:tidy-neighbours}
\begin{aligned}
\mathit{nb}_{\mathrm{l}}(p) &= \mathit{op}(B_t)^{\uparrow B_t}_{\mathrm{l}}
  \cdots \mathit{op}(B_h)^{\uparrow B_h}_{\mathrm{l}}
  && \text{if } p = \kw{gap}(\pi, 0) \text{ and } h > 0 \\
\mathit{nb}_{\mathrm{r}}(q) &= \mathit{cl}(B'_{h'})^{\downarrow
  B'_{h'}}_{\mathrm{r}} \cdots \mathit{cl}(B'_{t'})^{\downarrow
  B'_{t'}}_{\mathrm{r}}
  && \text{if } q = \kw{gap}(\pi', |\bar g'|),\ h' > 0, \text{ and }
     \mathit{cl}(B'_{h'}) \text{ is defined}
\end{aligned}
\end{equation}
and $\epsilon$ otherwise, where $\bar g'$ are the children of $B'_{h'}$;
$t \ge 1$ is least such that $B_l$ is the first child of $B_{l-1}$ for $t <
l \le h$, and $t' \ge 1$ least such that $B'_l$ is the last child of
$B'_{l-1}$ and $\mathit{cl}(B'_{l-1})$ is defined for $t' < l \le h'$.  An
opening of an inline block followed by a closing of it cancels:
\begin{equation}\label{eq:tidy-cancel}
\bar\omega\ \sigma^{\uparrow B}_{o}\ \sigma'^{\downarrow B}_{o'}\
  \bar\omega' \leadsto \bar\omega\ \bar\omega'
\qquad \text{if } \mathit{cls}(B) = \kw{inline}
\end{equation}
The items an edit writes also hold one \emph{caret marker} $\crt$, where the
caret goes.  The marker is never removed, and a pair cancels across it:
\begin{equation}\label{eq:tidy-caret}
\bar\omega\ \sigma^{\uparrow B}_{o}\ \crt\ \sigma'^{\downarrow B}_{o'}\
  \bar\omega' \leadsto \bar\omega\ \crt\ \bar\omega'
\qquad \text{if } \mathit{cls}(B) = \kw{inline}
\end{equation}
Two redexes never share an item, so $\leadsto$ has unique normal forms,
$\mathit{nf}(\bar\omega)$.  Then
\begin{equation}\label{eq:tidy}
\mathit{tidy}_m(a, a', \bar\omega) =
\begin{cases}
\kw{unchanged} & \text{if } a = a' \text{ and } \bar\omega = \crt \\
\mathit{splice}_m(a_\circ, a'_\circ, |\bar\omega_1|_{\mathrm{s}},
  |\bar\omega_2|_{\mathrm{s}}) &
  \text{otherwise, for } \bar\omega_1\,\crt\,\bar\omega_2 =
  \mathit{nf}(\mathit{drop}(\mathit{nb}_{\mathrm{l}}(p)\ \bar\omega\
  \mathit{nb}_{\mathrm{r}}(q)))
\end{cases}
\end{equation}
where $\mathit{drop}$ removes empty texts, $|\cdot|_{\mathrm{s}}$ keeps the
items of origin $\mathrm{s}$ and the unmarked ones and erases their marks,
and the splice is widened over the messages removed outside it.  If
$\mathit{nb}_{\mathrm{l}}(p) = \sigma_1 \cdots \sigma_e$ and $k$ of them
remain in $\bar\omega_\circ = \bar\omega_1\,\crt\,\bar\omega_2$, they are $\sigma_1 \cdots \sigma_k$, and
$a_\circ$ is $a$ if $k = e$ and $\mathit{before}(\sigma_{k+1})$ otherwise;
if $\mathit{nb}_{\mathrm{r}}(q) = \sigma'_1 \cdots \sigma'_{e'}$, innermost
first, and $k'$ of them remain, they are the last $k'$, and $a'_\circ$ is
$a'$ if $k' = e'$ and $\mathit{after}(\sigma'_{e'-k'})$ otherwise.  So
deleting all of an emphasis removes the emphasis rather than emptying it,
and Enter at the start of an emphasis leaves an empty paragraph but no
empty emphasis.  An edit that writes nothing at one anchor changes nothing,
so a caret inside an empty block that the document already has leaves it
alone; the caret too stays where it was.

\paragraph{Examples.}  Deleting all of the emphasis of $W_1$, from
$\kw{gap}(0.1, 0)$ to $\kw{gap}(0.1, 1)$, writes no items of its own.  The
first point is the first gap of $E$, which is not the first child of $P$,
and the second is the last gap of $E$, which is the last child of $P$:
\[
\mathit{nb}_{\mathrm{l}} = \str{emphasis}^{\uparrow E}_{\mathrm{l}}
\qquad
\mathit{nb}_{\mathrm{r}} = \str{endEmphasis}^{\downarrow E}_{\mathrm{r}}\
  \str{endParagraph}^{\downarrow P}_{\mathrm{r}}
\]
The first two cancel, so the splice is widened to the anchors
\[
\mathit{before}(\str{emphasis}) = \kw{txt}(n_1, 3)
\qquad
\mathit{after}(\str{endEmphasis}) = \kw{cut}(n_2, 1)
\]
and the source becomes
\[
\str{[\bs{}paragraph\bs]Hi\ [\bs{}endParagraph\bs]}
\]
[edit-delete-tidy].  Enter at the start of the emphasis of
\[
\str{[\bs{}paragraph.\ emphasis\bs]abc[\bs{}endEmphasis.\ endParagraph\bs]}
\]
writes the four statements $\str{endEmphasis}$, $\str{endParagraph}$,
$\str{paragraph}$, and $\str{emphasis}$ after the neighbours
$\str{paragraph}^{\uparrow}_{\mathrm{l}}$ and
$\str{emphasis}^{\uparrow}_{\mathrm{l}}$; the emphasis opened and closed
cancels, the paragraph, which is not inline, does not, and the splice is
widened to $\mathit{before}(\str{emphasis})$ [edit-split-tidy]:
\[
\begin{array}{@{}l@{}}
\str{[\bs{}paragraph\bs][\bs{}endParagraph.\ paragraph.\ emphasis\bs]} \\
\str{abc[\bs{}endEmphasis.\ endParagraph\bs]}
\end{array}
\]
The caret inside the empty emphasis of
\[
\str{[\bs{}paragraph\bs]a[\bs{}emphasis.\ endEmphasis\bs]b[\bs{}endParagraph\bs]}
\]
has both neighbours, which would
cancel, but an empty edit there changes nothing [edit-identity-empty].

\subsection{Edits}
\label{sec:edits}

The edits of the view are
\begin{equation}\label{eq:edits}
e ::= \kw{replace}(c, c', \bar\omega, \bar\omega') \mid \kw{split}(c, i) \mid
  \kw{style}(c, c', x) \mid \kw{setString}(\pi.\varphi_j, u) \mid
  \kw{detach}(\pi.\varphi_j)
\end{equation}
where $c$ and $c'$ are carets, and $p = \mathit{pt}(c)$ and $q =
\mathit{pt}(c')$ their points; a replacement writes $\bar\omega$, then the
caret, then $\bar\omega'$.  For a point $p$, let $B_0 \cdots B_h$ be the blocks on its path from the
root of its tree, $B_0$ the root.  For points $p \le q$ of one tree, with
blocks $B_0 \cdots B_h$ and $B'_0 \cdots B'_{h'}$, let $c$ be the depth of
the deepest block on both paths; then
\begin{equation}\label{eq:edit-blocks}
D = B_c \qquad
L_i = B_{c+i} \ \ (1 \le i \le a = h - c) \qquad
R_i = B'_{c+i} \ \ (1 \le i \le b = h' - c)
\end{equation}
are the deepest common block and the blocks below it on the path of $p$
and on the path of $q$, outermost first.  For a block $B$ with children
$g_0 \cdots g_{r-1}$, $\mathit{muts}_{[l, l')}(B)$ is the sequence of the
$\sigma$ with $g_i = \kw{mut}(\sigma)$ for $l \le i < l'$.  The mutating
messages of a block that lie between $p$ and $q$ are
\begin{equation}\label{eq:edit-muts}
M(L_i) = \mathit{muts}_{[\mathit{lo}(L_i), r)}(L_i)
\qquad
M(R_i) = \mathit{muts}_{[0, \mathit{hi}(R_i))}(R_i)
\qquad
M(D) = \mathit{muts}_{[\mathit{lo}(D), \mathit{hi}(D))}(D)
\end{equation}
where, for a block on the path of $p$ whose child $i'$ is the next on the
path, $\mathit{lo} = i' + 1$, and for the last block of the path,
$\mathit{lo} = i$ for $p = \kw{gap}(\pi, i)$ and $i + 1$ for $p =
\kw{in}(\pi, i, k)$; $\mathit{hi}$ is $i'$, or $i$ for both kinds of
point, on the path of $q$.  Two blocks \emph{join}, and the depth $j$ to
which the paths of $p$ and $q$ join is
\begin{equation}\label{eq:edit-join}
\begin{gathered}
B \sim B' \iff \mathit{merge}(\kappa_B) \wedge
  \lfloor \mathit{op}(B) \rfloor = \lfloor \mathit{op}(B') \rfloor
\\
j = \max\{\, j \le \min(a, b) \mid \forall i,\ 1 \le i \le j.\ L_i \sim R_i
  \wedge M_{i-1} = \epsilon \,\}
\end{gathered}
\end{equation}
for $\kappa_B$ the state of $B$, $M_0 = M(D)$, and $M_i = M(L_i)\,M(R_i)$
for $i > 0$.

\paragraph{Replacing.}
Replacing what lies between two carets by new items $\bar\omega\,\bar\omega'$
is the edit for typing, deleting, pasting, and inserting a block:
\begin{multline}\label{eq:edit-replace}
\mathit{replace}(c, c', \bar\omega, \bar\omega') =
\mathit{tidy}_m(\mathit{anchor}(c), \mathit{anchor}(c'),\
  \bar\omega\,\crt\,\bar\omega'\ M(L_a)\,\mathit{cl}(L_a) \cdots M(L_{j+1})\,\mathit{cl}(L_{j+1}) \\
  M_j\ \mathit{op}(R_{j+1})\,M(R_{j+1}) \cdots \mathit{op}(R_b)\,M(R_b))
\end{multline}
where $m = \mathit{mk}(p) = \mathit{mk}(q)$.  Everything shown between $p$
and $q$ goes.  A block that holds $p$ but not $q$ keeps what lies before
$p$ and is closed after $\bar\omega$; one that holds $q$ but not $p$ keeps
what lies after $q$ and is opened again; and the outermost pairs of such
blocks that join are merged into one.  A mutating message renders nothing
where it lies, so it is kept and sent to the same block, and $j$ is lowered
where it could not be.  The keys of the view are replacements:
\begin{equation}\label{eq:edit-keys}
\begin{aligned}
\mathit{type}(c, u) &= \kw{replace}(c, c, \kw{text}(u), \epsilon) \\
\mathit{delete}(c, c') &= \kw{replace}(c, c', \epsilon, \epsilon) \\
\mathit{backspace}(\pi.i) &= \kw{replace}(\kw{gap}(\pi.(i{-}1),
  |\bar g|),\ \kw{gap}(\pi.i, 0),\ \epsilon, \epsilon)
  \qquad \text{for } i > 0
\end{aligned}
\end{equation}
where $\bar g$ are the children of $\tau@\pi.(i{-}1)$: Backspace at the
start of the block at $\pi.i$, and Delete at the end of its predecessor,
deletes from the end of the predecessor to the start of the block, which
merges the two when they join.  Inserting a block is replacing with its
template (Section~\ref{sec:lib-field}).

\paragraph{Examples.}
\begin{itemize}
\item Typing $\str{!}$ at the end of $\str{all}$ in $W_1$,
      $\mathit{type}(\kw{gap}(0.1, 1), \str{!})$, has $D = E$ and no $L_i$
      or $R_i$; it splices $\kw{text}(\str{!})$ at $\kw{txt}(n_2, 3)$, and
      the contents of $n_2$ become $\str{all!}$ [edit-type].  Replacing
      nothing there by nothing changes nothing [edit-identity].
\item In $\str{[\bs{}paragraph\bs]a\ [\bs{}emphasis\bs]bc[\bs{}endEmphasis\bs]\
      d[\bs{}endParagraph\bs]}$, typing $\str{X}$ over the selection from
      $\kw{in}(0.1, 0, 1)$, between $\str{b}$ and $\str{c}$, to
      $\kw{in}(0, 2, 1)$, before $\str{d}$, has $D$ the paragraph, $L_1$
      the emphasis, and $b = 0$.  The emphasis keeps $\str{b}$ and is
      closed after $\str{X}$ [edit-replace-across]:
      $\str{[\bs{}paragraph\bs]a\ [\bs{}emphasis\bs]bX[\bs{}endEmphasis\bs]d[\bs{}endParagraph\bs]}$.
\item Deleting from $\kw{in}(0, 0, 1)$ to $\kw{in}(1, 0, 1)$ in two
      paragraphs $\str{ab}$ and $\str{cd}$ has $D$ the root and $L_1$ and
      $R_1$ the paragraphs, which join, so $j = 1$ and neither
      $\mathit{cl}(L_1)$ nor $\mathit{op}(R_1)$ is written:
      $\str{[\bs{}paragraph\bs]ad[\bs{}endParagraph\bs]}$
      [edit-delete-join].  With $\str{title:\ [\bs{}T\bs]}$ sent to the
      root between them, $M(D)$ is not empty, so $j = 0$; the message is
      kept, and so are both paragraphs [edit-delete-keeps-title]:
      \[
      \str{[\bs{}paragraph\bs]a[\bs{}endParagraph.\ title:\
      [\bs{}T\bs].\ paragraph\bs]d[\bs{}endParagraph\bs]}
      \]
\end{itemize}

\paragraph{Splitting.}
Splitting the block $B_i$, $i \ge 1$, at the caret $c$ closes the blocks
from the innermost out to $B_i$ and opens them again, and puts the caret at
the start of the second half:
\begin{equation}\label{eq:edit-split}
\mathit{split}(c, B_i) = \mathit{tidy}_m(\mathit{anchor}(c),
  \mathit{anchor}(c),\ \mathit{cl}(B_h) \cdots \mathit{cl}(B_i)\
  \mathit{op}(B_i) \cdots \mathit{op}(B_h)\ \crt)
\end{equation}
which needs each of $B_i, \ldots, B_h$ to have a closing message.  The
second half receives none of the mutating messages of the first, so it
starts with the fields its creating message gave it.  Enter in a paragraph
splits the paragraph, and in a list item the item
(\ref{eq:lib-enter}).

\paragraph{Examples.}  Enter between $\str{Hi}$ and the space of $W_1$ is
$\kw{split}(\kw{in}(0, 0, 2), 1)$ [ex-enter], which gives
\[
\begin{array}{@{}l@{}}
\str{[\bs{}paragraph\bs]Hi[\bs{}endParagraph.\ paragraph\bs]\ } \\
\str{[\bs{}emphasis\bs]all[\bs{}endEmphasis.\ endParagraph\bs]}
\end{array}
\]
and the second paragraph's leading space is trimmed when it closes
[edit-split, edit-split-tree].  Backspace at the start of the second
paragraph is $\mathit{backspace}(\epsilon.1) = \kw{replace}(\kw{gap}(0,
1), \kw{gap}(1, 0), \epsilon)$ [ex-backspace]: the paragraphs join, so
neither closer nor opener is kept, and the splice gives back $W_1$
[edit-join].  Enter inside the emphasis, at
$\kw{in}(0.1, 0, 1)$, splits the paragraph too, $B_1$ and not $B_2$
[ex-enter-emph], so the emphasis is closed and opened again around the
break [ex-split-emph]:
\[
\begin{array}{@{}l@{}}
\str{[\bs{}paragraph\bs]Hi\ [\bs{}emphasis\bs]a} \\
\str{[\bs{}endEmphasis.\ endParagraph.\ paragraph.\ emphasis\bs]} \\
\str{ll[\bs{}endEmphasis.\ endParagraph\bs]}
\end{array}
\]

\paragraph{Styling.}
Toggling a style $x$ on the range from $p$ to $q$, $p < q$, needs
\begin{equation}\label{eq:style-ok}
\mathit{styleable}(p, q) \iff \mathit{cls}(D) \in \{\kw{para}, \kw{cell},
  \kw{inline}\} \wedge \forall i.\ \mathit{cls}(L_i) = \mathit{cls}(R_i) =
  \kw{inline}
\end{equation}
If some inline block on both paths was opened as $x$ is, the range is
styled and the style is removed from it.  The innermost such block is
$B = B_l$ for
\[
l = \max\{\, l \le c \mid \mathit{cls}(B_l) = \kw{inline} \wedge
  \lfloor \mathit{op}(B_l) \rfloor = \lfloor \mathit{open}(x) \rfloor \,\}
\]
and, with $C_1 \cdots C_e = B_{l+1} \cdots B_c$ the blocks below it down to
$D$,
\begin{equation}\label{eq:edit-unstyle}
\begin{aligned}
\bar\omega_p &= \mathit{cl}(L_a) \cdots \mathit{cl}(L_1)\
  \mathit{cl}(C_e) \cdots \mathit{cl}(C_1)\ \mathit{cl}(B)\
  \mathit{op}(C_1) \cdots \mathit{op}(C_e)\
  \mathit{op}(L_1) \cdots \mathit{op}(L_a) \\
\bar\omega_q &= \mathit{cl}(R_b) \cdots \mathit{cl}(R_1)\
  \mathit{cl}(C_e) \cdots \mathit{cl}(C_1)\ \mathit{op}(B)\
  \mathit{op}(C_1) \cdots \mathit{op}(C_e)\
  \mathit{op}(R_1) \cdots \mathit{op}(R_b)
\end{aligned}
\end{equation}
Otherwise the style is applied:
\begin{equation}\label{eq:edit-restyle}
\begin{gathered}
\bar\omega_p = \mathit{cl}(L_a) \cdots \mathit{cl}(L_1)\ \mathit{open}(x)\
  \mathit{op}(L_1) \cdots \mathit{op}(L_a)
\\
\bar\omega_q = \mathit{cl}(R_b) \cdots \mathit{cl}(R_1)\ \mathit{close}(x)\
  \mathit{op}(R_1) \cdots \mathit{op}(R_b)
\end{gathered}
\end{equation}
and in both cases
\begin{equation}\label{eq:edit-style}
\mathit{style}(c, c', x) = \mathit{tidy}_m(\mathit{anchor}(c),
  \mathit{anchor}(c'),\ \bar\omega_p\ \mathit{items}(\mathit{anchor}(c),
  \mathit{anchor}(c'))\ \crt\ \bar\omega_q)
\end{equation}
so the caret ends at the end of the range, in the blocks its text now lies
in.
The blocks open at $p$ are closed there and opened again inside the new
style, and likewise at $q$, so the style nests properly whatever the range
crosses.  A range only partly styled is styled throughout, so toggling
twice gives back the source's tree only when the range was plain or
styled throughout.

\paragraph{Examples.}
\begin{itemize}
\item Toggling emphasis on $\str{Hi}$ in $W_1$, from $\kw{gap}(0, 0)$ to
      $\kw{in}(0, 0, 2)$, finds no emphasis on both paths and applies it
      [edit-style-apply]:
      \[
      \begin{array}{@{}l@{}}
      \str{[\bs{}paragraph\bs][\bs{}emphasis\bs]Hi[\bs{}endEmphasis\bs]\ } \\
      \str{[\bs{}emphasis\bs]all[\bs{}endEmphasis.\ endParagraph\bs]}
      \end{array}
      \]
\item Toggling emphasis on $\str{all}$, from $\kw{gap}(0.1, 0)$ to
      $\kw{gap}(0.1, 1)$, finds $B = E$ and removes it; the statements
      $\str{endEmphasis}\ \str{emphasis}$ that it writes at the two ends
      cancel against $E$'s own, and the source becomes
      $\str{[\bs{}paragraph\bs]Hi\ all[\bs{}endParagraph\bs]}$
      [edit-style-remove].
\item In the source of [edit-replace-across] above, making the range from inside the
      emphasis to before $\str{d}$ strong closes the emphasis at $p$ and
      opens it again inside the strong text [edit-style-across]:
      \[
      \begin{array}{@{}l@{}}
      \str{[\bs{}paragraph\bs]a\ [\bs{}emphasis\bs]b[\bs{}endEmphasis.\ strong.\ emphasis\bs]} \\
      \str{c[\bs{}endEmphasis\bs]\ [\bs{}endStrong\bs]d[\bs{}endParagraph\bs]}
      \end{array}
      \]
\end{itemize}

\paragraph{Strings and detaching.}
For a field $\varphi = \langle v, \sigma, k \rangle$ of the block at $\pi$,
whose message $\sigma$ has the arguments $\bar\eta$, let $W = I(d)$ for $d
= \mathit{src}(\mathit{mk}(\pi))$, the source that holds $\sigma$.
Setting the string of $\varphi$ to $u$ and detaching $\varphi$ rewrite the
text of its argument:
\begin{equation}\label{eq:edit-string}
\begin{aligned}
\mathit{setString}(\varphi, u) &= W[0, i)\ \mathit{enc}(u)\ W[i', |W|)
  && \text{for } \mathit{str}(\eta_k) = [i, i') \\
\mathit{detach}(\varphi) &= W[0, i)\ \mathcal{Q}_{\mathrm{v}}(v)\
  W[i', |W|)
  && \text{for } \mathit{span}(\eta_k) = [i, i')
\end{aligned}
\end{equation}
so that after a detach the value is written in the source as a literal.  A
field can be detached when its message lies in a markup that is not
computed and the host can write its value, and what the view offers for a
field is
\begin{equation}\label{eq:detachable}
\begin{gathered}
\mathit{detachable}(\varphi) \iff \mathit{prov}(\mathit{mk}(\pi)) \neq
  \kw{computed} \wedge \mathcal{Q}_{\mathrm{v}}(v) \text{ is defined}
\\
\mathit{offers}(\varphi) =
\begin{cases}
\kw{detach}(\pi.\varphi_j)\ \kw{giveUp} & \text{if }
  \mathit{detachable}(\varphi) \\
\kw{giveUp} & \text{otherwise}
\end{cases}
\end{gathered}
\end{equation}
for $\varphi$ the $j$th field of the block at $\pi$.

\subsection{The edit relation}

\begin{equation}\label{eq:edit-results}
\varrho ::= \kw{edited}(d, W', \hat c) \mid \kw{refused}(\bar o)
\qquad
\hat c ::= k \mid \kw{same}
\qquad
o ::= \kw{detach}(\pi.\varphi_j) \mid \kw{giveUp}
\end{equation}
An edit either gives the new text $W'$ of one source $d$, with the offset
$k$ in it of the caret or $\kw{same}$ when the caret stays where it was, or
is refused with the offers the view should make.  A replacement, split, or
style gives
\[
\mathit{result}_d(\langle W', k \rangle) = \kw{edited}(d, W', k)
\qquad
\mathit{result}_d(\kw{unchanged}) = \kw{edited}(d, I(d), \kw{same})
\]  The edit relation $I \vdash e
\Rightarrow \varrho$ is relative to the host state $\chi$ and root state
$\kappa_0$ with which the document is interpreted:
\begin{gather}
\rulelabel{Document}\quad
\frac{\mathit{parse}(I(d_0)) = \kw{parsed}(\bar n)
      \qquad \langle \chi, \kappa_0, \bar n \rangle \Downarrow_{\mathrm{I}}
        \kw{done}(\chi', \tau)}
     {\mathit{doc}(I) = \langle \bar n, \tau \rangle}
\end{gather}
In the rules below, $\mathit{doc}(I) = \langle \bar n, \tau \rangle$, the
points and paths are those of $\tau$, $m = \mathit{mk}(p)$ for an edit with
the points $p$ and $q$, and $p \asymp q$ holds when $p$ and $q$ lie in one
tree.  For $\kw{setString}$ and $\kw{detach}$, $\varphi = \langle v,
\sigma, k \rangle$ is the $j$th field of $\tau@\pi$, $m =
\mathit{mk}(\pi)$, and $\eta_k$ is the $k$th argument of $\sigma$.
\begin{gather}
\rulelabel{Edit-Replace}\quad
\frac{p \asymp q \qquad \mathit{prov}(m) \neq \kw{computed}}
     {I \vdash \kw{replace}(c, c', \bar\omega, \bar\omega') \Rightarrow
      \mathit{result}_{\mathit{src}(m)}(\mathit{replace}(c, c', \bar\omega,
        \bar\omega'))}
\\
\rulelabel{Edit-Split}\quad
\frac{\mathit{prov}(m) \neq \kw{computed} \qquad 1 \le i \le h \qquad
      \mathit{cl}(B_l) \text{ is defined for } i \le l \le h}
     {I \vdash \kw{split}(c, i) \Rightarrow
      \mathit{result}_{\mathit{src}(m)}(\mathit{split}(c, B_i))}
\\
\rulelabel{Edit-Style}\quad
\frac{p \asymp q \qquad \mathit{prov}(m) \neq \kw{computed} \qquad
      \mathit{styleable}(p, q)}
     {I \vdash \kw{style}(c, c', x) \Rightarrow
      \mathit{result}_{\mathit{src}(m)}(\mathit{style}(c, c', x))}
\\
\rulelabel{Edit-Set-String}\quad
\frac{\mathit{prov}(m) \neq \kw{computed} \qquad
      \mathit{str}(\eta_k) \text{ is defined}}
     {I \vdash \kw{setString}(\pi.\varphi_j, u) \Rightarrow
      \kw{edited}(\mathit{src}(m), \mathit{setString}(\varphi, u),
        \kw{same})}
\\
\rulelabel{Edit-Detach}\quad
\frac{\mathit{detachable}(\varphi)}
     {I \vdash \kw{detach}(\pi.\varphi_j) \Rightarrow
      \kw{edited}(\mathit{src}(m), \mathit{detach}(\varphi), \kw{same})}
\end{gather}
An edit aimed at computed text is offered, as the design notes ask, a
choice: to detach, where that is possible, or to give up.  For a point $p$,
$\varphi(p)$ is the field its path entered last, and $\mathit{offers}$ of
no field is $\kw{giveUp}$.  In Refuse-Computed, $e$ is a $\kw{replace}$, a
$\kw{split}$, or a $\kw{style}$ whose first point is $p$, and $p \asymp q$
when it has a second point $q$.
\begin{gather}
\rulelabel{Refuse-Trees}\quad
\frac{\neg(p \asymp q)}
     {I \vdash e \Rightarrow \kw{refused}(\kw{giveUp})}
  \quad e \in \{\kw{replace}(c, c', \bar\omega, \bar\omega'),
    \kw{style}(c, c', x)\}
\\
\rulelabel{Refuse-Computed}\quad
\frac{\mathit{prov}(\mathit{mk}(p)) = \kw{computed}}
     {I \vdash e \Rightarrow \kw{refused}(\mathit{offers}(\varphi(p)))}
\\
\rulelabel{Refuse-Split}\quad
\frac{\mathit{prov}(m) \neq \kw{computed} \qquad
      \neg(1 \le i \le h \wedge \mathit{cl}(B_l) \text{ is defined for } i
        \le l \le h)}
     {I \vdash \kw{split}(c, i) \Rightarrow \kw{refused}(\kw{giveUp})}
\\
\rulelabel{Refuse-Style}\quad
\frac{p \asymp q \qquad \mathit{prov}(m) \neq \kw{computed} \qquad
      \neg\,\mathit{styleable}(p, q)}
     {I \vdash \kw{style}(c, c', x) \Rightarrow \kw{refused}(\kw{giveUp})}
\\
\rulelabel{Refuse-String}\quad
\frac{\neg(\mathit{prov}(m) \neq \kw{computed} \wedge
        \mathit{str}(\eta_k) \text{ is defined})}
     {I \vdash \kw{setString}(\pi.\varphi_j, u) \Rightarrow
      \kw{refused}(\mathit{offers}(\varphi))}
\\
\rulelabel{Refuse-Detach}\quad
\frac{\neg\,\mathit{detachable}(\varphi)}
     {I \vdash \kw{detach}(\pi.\varphi_j) \Rightarrow
      \kw{refused}(\kw{giveUp})}
\end{gather}
The premises of the rules for each kind of edit are exclusive and
exhaustive, so every edit has exactly one result.  After a detach, the
view repeats the edit on the reinterpreted document, where the field is now
local.  An edit of a template is not refused: it goes to the method's
source, as the design notes require, and the view shows its scope, since
every document that uses the template changes with it.  Detaching a
template's field instead gives the document its own copy.

\paragraph{Examples.}
\begin{itemize}
\item In $\str{[\bs{}title:\ actTitle\bs]}$ the title is a string that the
      host method $\str{actTitle}$ computes.  Setting it is refused with
      $\kw{detach}(\varphi_0)\ \kw{giveUp}$ [edit-string-computed], and
      detaching it writes its value as a string literal,
      $\str{[\bs{}title:\ 'Act\ I'\bs]}$, after which it can be set
      [edit-detach].
\item Detaching the title of $W_{\mathrm{t}}$ puts a copy of the
      template's literal in the document [edit-detach-template]:
      \[
      \str{[\bs{}title:\ [\bs{}The\ Play\bs]\bs]}
      \]
\item Typing in the computed title of $W_{\mathrm{c}}$ is refused with
      $\kw{giveUp}$ alone: the field cannot be detached, since the host
      cannot write a tree whose literal has no source in the image
      [edit-computed-tree].
\item Deleting from the title of $W_{\mathrm{l}}$ into its paragraph
      [edit-two-trees], Enter in a paragraph that is never closed
      [edit-split-unclosed], and styling a range across two paragraphs
      [edit-style-flow] are refused with $\kw{giveUp}$.
\end{itemize}


\paragraph{The caret after an edit.}  For $\kw{edited}(d, W', k)$, let $I'
= I[d := W']$ and $\mathit{doc}(I') = \langle \bar n', \tau' \rangle$.  The
edited markup is at the same path $\pi_m$ (the edit changes nothing outside
it, by locality), and the caret is placed in it:
\begin{gather}
\rulelabel{Caret}\quad
\frac{\mathit{doc}(I') = \langle \bar n', \tau' \rangle \qquad
      m' = \mathit{mk}(\pi_m) \text{ in } \tau' \qquad
      a = \mathit{at}_{m'}(k)}
     {I' \vdash \mathit{caret}(\pi_m, k) = \kw{at}(\pi_m.\mathit{place}(a), a)}
\end{gather}
For $\kw{same}$, the caret is the one the edit was given.

\paragraph{Examples.}
\begin{itemize}
\item Typing a space at the end of $\str{[\bs{}paragraph\bs]Hello[\bs{}endParagraph\bs]}$,
      at $\kw{gap}(0, 1)$, writes the space and puts the caret after it, at
      offset $19$ [caret-space].  The space is trimmed, so the caret is
      $\kw{at}(\kw{gap}(0, 1), \kw{txt}(n_1, 6))$ [caret-space-caret], and
      typing $\str{w}$ there gives $\str{Hello}\vs\str{w}$
      [caret-after-space].  At the exact caret, whose anchor is
      $\kw{txt}(n_1, 5)$, it would give $\str{Hellow}\vs$
      [caret-exact-trims].
\item A second space typed after the one of
      $\str{[\bs{}paragraph\bs]a}\vs\str{b[\bs{}endParagraph\bs]}$
      collapses into it; the caret is $\kw{at}(\kw{in}(0, 0, 2),
      \kw{txt}(n_1, 3))$, still before the $\str{b}$ [caret-collapsed].
\item Enter between $\str{Hi}$ and the space of $W_1$ leaves the caret at
      the start of the second paragraph, $\kw{at}(\kw{gap}(1, 0),
      \kw{txt}(n_2, 0))$, the space being trimmed [caret-enter].  Deleting
      all of the emphasis leaves it where the emphasis was,
      $\kw{at}(\kw{gap}(0, 1), \kw{txt}(n_1, 3))$: the emphasis's opener and
      closer cancel across the marker [caret-delete-emph].
\item Typing $\str{-}$ in the title of $\str{[\bs{}title:\ [\bs{}Hamlet\bs]\bs]}$
      places the caret in the literal, below the field: $\kw{at}(\kw{in}(\varphi_0, 0,
      4), \kw{txt}(n_0, 4))$ [caret-title].
\end{itemize}

\subsection{Laws}

\paragraph{Locality.}
By construction an edit changes only the source of its points, and in it
only the text between the bounds of its anchors as tidying widened them:
if $I \vdash e \Rightarrow \kw{edited}(d, W', \hat c)$ for a $\kw{replace}$,
$\kw{split}$, or $\kw{style}$ with the widened anchors $a_\circ, a'_\circ$
of $m$, then $d = \mathit{src}(m)$ and
\begin{equation}\label{eq:law-locality}
W' = I(d)[0, \mathit{L}(a_\circ))\ w\ I(d)[\mathit{R}(a'_\circ), |I(d)|)
\end{equation}
for some $w$; for $\kw{setString}$ and $\kw{detach}$ the same holds of the
argument's span.  A local edit never changes a template, and an edit of a
template never changes the document.

\paragraph{Identity.}
\begin{equation}\label{eq:law-identity}
I \vdash \kw{replace}(c, c, \epsilon, \epsilon) \Rightarrow
  \kw{edited}(\mathit{src}(m), I(\mathit{src}(m)), \kw{same})
\quad\text{when } \mathit{prov}(m) \neq \kw{computed}
\end{equation}
by the first case of (\ref{eq:tidy}).

\paragraph{Typing (obligation).}
If the anchors of $p$ and $q$ are $\kw{txt}(n, o)$ and $\kw{txt}(n, o')$
in the contents $u$ of one node $n$, whose text is the child $g$ of the
block $B$ at the end of the path of $p$, then typing replaces $g$ alone:
\begin{equation}\label{eq:edit-typing}
\mathit{replace}(p, q, \kw{text}(u''), \epsilon) \text{ gives } W', \text{ whose
tree is } \approx\ \tau[g := g']
\end{equation}
where $g'$ is the child that the contents $u[0, o)\,u''\,u[o', |u|)$ give in
$B$ (\ref{eq:text-child}), trimmed as closing $B$ trims it; if that is no
child, $g$ is removed.  The statements of the source are unchanged, so by
(H7) every message has the effect it had before.

\paragraph{Library edits (obligations).}  For the standard library, with
$\tau'$ the tree of the result:
\begin{itemize}
\item Backspace at the start of the second of two consecutive paragraphs
      $P_1 P_2$ gives one paragraph holding the children of both,
      $\tau' \approx \tau[P_1 P_2 := P]$ with $P$ the paragraph of
      $P_1$ whose children are those of $P_1$ then those of $P_2$, trimmed.
\item Splitting a paragraph $P$ at a point gives two paragraphs, $\tau'
      \approx \tau[P := P_1 P_2]$ with the children before and after the
      point, each trimmed.
\item Inserting the template of $x$ at a point whose block understands $x$
      adds no error: $\mathit{errors}(\tau') = \mathit{errors}(\tau)$.
\item Toggling a style twice on a range of plain text gives back the tree:
      $\tau'' \approx \tau$.
\end{itemize}
The regressions check instances of each [edit-join, edit-split-tree,
edit-insert-block-tree, edit-style-apply, edit-style-remove].

\subsection{The standard library's field interface}
\label{sec:lib-field}

The fields of the standard library's blocks are the values it displays:
\begin{equation}\label{eq:lib-fields}
\begin{array}{@{}ll@{}}
\kw{doc}(t, t', \bar a) & t, t', \bar a \text{ from } \str{title:},
  \str{subtitle:}, \str{author:} \\
\kw{section}(d, h, \ell),\ \kw{quote}(a),\ \kw{figure}(c, \ell)
  & h, a, c \text{ from } \str{section:}, \str{attribution:},
  \str{caption:} \\
\kw{verbatim}(u),\ \kw{codeBlock}(u, x),\ \kw{html}(u),\ \kw{link}(u)
  & u \text{ from the creating message} \\[-1ex]
\end{array}
\end{equation}
each the first argument of its message; a field that is $\kw{nil}$ is not
shown, and labels and code languages are not displayed.  The blocks that
merge are paragraphs, the three styles, links (with the same target, since
their creating messages must be the same), lists, numbered lists, and items;
sections, tables, rows, cells, quotes, and figures do not, so deleting a
selection that spans two cells clears the cells but keeps both.  The styles
are $\str{emphasis}$, $\str{strong}$, and $\str{code}$, closed by
$\str{endEmphasis}$, $\str{endStrong}$, and $\str{endCode}$.

The \emph{template} of an opener $x$ is the opener followed by its closer,
with the view's caret between them, and a table's holds one row of one
cell; it is inserted only where the block at the point understands $x$:
\begin{equation}\label{eq:lib-templates}
\begin{gathered}
\mathit{tmpl}(x) = x\ \crt\ \mathit{closer}(x) \quad (x \neq \str{table})
\\
\mathit{tmpl}(\str{table}) = \str{table}\ \str{row}\ \str{cell}\ \crt\
  \str{endCell}\ \str{endRow}\ \str{endTable}
\\
\mathit{insert}(c, x) = \kw{replace}(c, c, \bar\omega, \bar\omega') \quad
  \text{for } \mathit{tmpl}(x) = \bar\omega\,\crt\,\bar\omega', \text{ if }
  x \in \mathit{model}(\kappa_p)
\end{gathered}
\end{equation}
for $\kappa_p$ the state of the block of $p = \mathit{pt}(c)$, so the
caret lands inside the new block.  Here $\mathit{closer}$
gives $\str{endParagraph}$ for $\str{paragraph}$, $\str{endRow}$ for
$\str{row}$ and $\str{headerRow}$, and so on for every
opener of (\ref{eq:lib-openers}).  The template of a display is the
message with an empty string, as $\str{verbatim:\ ''}$, and that of an
inline-argument message the message with an empty literal, as
$\str{title:\ [\bs\bs]}$, which is how a reader gives a document a title it
lacks.  Enter splits the innermost paragraph or list item:
\begin{equation}\label{eq:lib-enter}
\mathit{enter}(c) = \kw{split}(c, l) \quad \text{for } l = \max\{\, l \ge 1
  \mid B_l \text{ is a } \kw{para} \text{ or } \kw{item} \text{ block}
  \,\}
\end{equation}

\paragraph{Examples.}  At the gap between the paragraphs of
$W_{\mathrm{p}}$ the root understands $\str{paragraph}$ but not
$\str{row}$, and at the start of a paragraph the paragraph understands
$\str{emphasis}$ but not $\str{paragraph}$ [ex-insertable,
ex-insertable-row, ex-insertable-emph, ex-insertable-para].  Inserting a
paragraph there writes its opener, the caret, and its closer [ex-insert];
the bracket of the new paragraph is split at the caret [caret-template,
caret-template-tree],
\[
\begin{array}{@{}l@{}}
\str{[\bs{}paragraph\bs]ab[\bs{}endParagraph\bs][\bs{}paragraph\bs]} \\
\str{[\bs{}endParagraph\bs][\bs{}paragraph\bs]cd[\bs{}endParagraph\bs]}
\end{array}
\]
and the caret is $\kw{at}(\kw{gap}(1, 0), \kw{txt}(n_3, 0))$, inside the
new paragraph [caret-template-caret].  A table's template puts the caret in
its one cell [ex-insert-table].
Writing the two statements of a paragraph's template with the caret after
both, after the last block of
$\str{[\bs{}paragraph\bs]ab[\bs{}endParagraph\bs]}$, adds an empty
paragraph in one bracket, $\str{[\bs{}paragraph.\ endParagraph\bs]}$
[edit-insert-block, edit-insert-block-tree].

\end{document}
